\documentclass[11pt,twocolumn,a4paper]{article}

\usepackage{amsmath,amssymb,amsthm}
\usepackage{mathtools}
\usepackage{physics}
\usepackage{geometry}
\usepackage{hyperref}
\usepackage{cleveref}
\usepackage{graphicx}
\usepackage{xcolor}
\usepackage{booktabs}
\usepackage{array}
\usepackage{enumitem}
\usepackage{natbib}
\usepackage{microtype}
\usepackage{bm}
\usepackage{bbm}

\geometry{margin=2.5cm}

\theoremstyle{plain}
\newtheorem{theorem}{Theorem}[section]
\newtheorem{lemma}[theorem]{Lemma}
\newtheorem{proposition}[theorem]{Proposition}
\newtheorem{corollary}[theorem]{Corollary}

\theoremstyle{definition}
\newtheorem{definition}[theorem]{Definition}
\newtheorem{axiom}[theorem]{Axiom}
\newtheorem{example}[theorem]{Example}
\newtheorem{remark}[theorem]{Remark}

\theoremstyle{remark}
\newtheorem{note}[theorem]{Note}

\DeclareMathOperator{\tr}{tr}
\DeclareMathOperator{\rank}{rank}
\DeclareMathOperator{\Span}{span}
\DeclareMathOperator{\vol}{vol}
\DeclareMathOperator{\diam}{diam}
\newcommand{\R}{\mathbb{R}}
\newcommand{\C}{\mathbb{C}}
\newcommand{\N}{\mathbb{N}}
\newcommand{\Z}{\mathbb{Z}}
\newcommand{\calH}{\mathcal{H}}
\newcommand{\calB}{\mathcal{B}}
\newcommand{\calF}{\mathcal{F}}
\newcommand{\calP}{\mathcal{P}}
\newcommand{\calT}{\mathcal{T}}
\newcommand{\Omega}{\Omega}
\newcommand{\hbar}{\hbar}
\newcommand{\half}{\tfrac{1}{2}}

\hypersetup{
  colorlinks=true,
  linkcolor=blue!60!black,
  citecolor=green!50!black,
  urlcolor=blue!70!black
}

\title{\textbf{The Resonant Spectral Stack}\\[0.5em]
\large Stacked-Fluid Liquid Mirrors, Loop-Coupled Temporal Echoes,\\
and the Triple Observation Identity as a Unified Instrument Architecture}

\author{
Kundai Farai Sachikonye\\[0.3em]
AIMe Registry for Artificial Intelligence\\
Experimental Bioinformatics, Maximus-von-Imhof-Forum 3\\
85354 Freising, Germany\\
\texttt{kundai.sachikonye@bitspark.com}
}

\date{2026}

\begin{document}

\maketitle

\begin{abstract}
We construct a unified optical instrument — the \emph{resonant spectral stack} — by synthesising three independently motivated physical structures: (i) a stacked-fluid liquid-mirror telescope in which $N$ superposed fluid layers of differing Cauchy-dispersion coefficients act as a wavelength-indexed refraction network; (ii) a loop-coupling cavity in which multiply-reflected rays accumulate wavelength-dependent temporal offsets that are formally equivalent to the Shapiro gravitational delay; and (iii) the \emph{triple observation identity}, which asserts that a single ray traversal simultaneously realises Beer--Lambert extinction, chromatographic retention, and Kirchhoff circuit balance as three coordinate projections of one underlying partition operation.

The mathematical foundation is laid in Part~I by deriving, from Poincar\'{e} recurrence and the Birkhoff ergodic theorem, that every bounded measure-preserving system necessarily supports a discrete frequency spectrum, hierarchical partition coordinates $(n,l,m,s)$ with ranges fixed by $\mathrm{SO}(3)$ geometry and the topology $\pi_1(\mathrm{SO}(3))=\mathbb{Z}_2$, a shell-capacity formula $C(n)=2n^2$, and a partition Lagrangian whose equations of motion reproduce the Lorentz force without postulating the electromagnetic field.

Part~II introduces the stacked-fluid architecture. The key result is the \emph{Spectral Separability Theorem}: if the fluid-layer matrix $A_\lambda$ has full column rank, the wavelength vector can be uniquely recovered from the observed output angles, making the stack a lossless hyperspectral spectrometer. Part~III derives the temporal-echo spectrum from the loop-coupling cavity, shows that the differential delay $\Delta\tau_n(\lambda)=L_n(\lambda)/c$ between spectral channels is exactly the Shapiro excess delay for an equivalent gravitational lens depth, and constructs the joint temporal--spectral distribution $\mathcal{S}(\lambda,\tau)$.

Part~IV proves the triple observation identity as a theorem: the Beer--Lambert law, the plate-theory chromatographic retention factor, and Kirchhoff's voltage law are identical up to change of variable. Part~V introduces the composition-inflation operator $T(n,d)=d(d{+}1)^{n-1}$ counting distinguishable trajectory histories, the three-component residue decomposition $A+R+P=\text{const}$, and the residue fraction $f_R=(b^d-1)/b^d$ which equals $26/27$ for $b=d=3$.

Part~VI assembles the unified instrument: a four-dimensional output tensor $\mathcal{O}[\lambda,n_\ell,z,\text{obs}]$ read off in five GPU passes, from which six independent derived instruments are extracted without additional hardware: three convergent routes to Newton's constant $G$; a composition-inflation visualiser; a computer-as-spectrometer that maps CPU clock, memory bus phase, LED emission frequency, and refresh polarity to $(n,l,m,s)$; a transport-extinction spectrometer; a zero-work categorical cooler in which temperature decreases as partition depth grows at fixed energy; and a sub-Planckian categorical clock achieving $\delta t_\text{cat}\sim 10^{-50}$ s from a 1950-mode log-spaced oscillator network.

Part~VII lists seven falsifiable experimental predictions distinguishing the resonant stack from standard thin-film or Fabry--P\'{e}rot spectrometry, with quantitative thresholds.
\end{abstract}

\tableofcontents
\newpage

%%%%%%%%%%%%%%%%%%%%%%%%%%%%%%%%%%%%%%%%%%%%%%%%%%%%%%%%%%
\part{Mathematical Foundations}
%%%%%%%%%%%%%%%%%%%%%%%%%%%%%%%%%%%%%%%%%%%%%%%%%%%%%%%%%%

\section{Bounded Measure-Preserving Systems and Poincar\'{e} Recurrence}
\label{sec:recurrence}

We begin with the minimal hypotheses sufficient to derive a discrete frequency spectrum from first principles, without assuming the Schr\"{o}dinger equation or the electromagnetic field as primitives.

\begin{axiom}[Bounded Phase Space]
\label{ax:bounded}
A physical system is a triple $(\Omega,\mu,\phi_t)$ where $\Omega\subset\R^{2N}$ is a phase-space region satisfying $0<\mu(\Omega)<\infty$, $\mu$ is a Borel measure on $\Omega$ preserved by the flow, i.e.\ $\mu(\phi_t(A))=\mu(A)$ for every measurable $A\subseteq\Omega$ and every $t\in\R$, and $\phi_t:\Omega\to\Omega$ is a one-parameter group of measurable bijections.
\end{axiom}

The bound $\mu(\Omega)<\infty$ is the sole physical content of the axiom; it is satisfied by any system whose accessible phase-space volume is finite — a condition that holds for every bound state encountered in laboratory physics.

\begin{theorem}[Poincar\'{e} Recurrence \cite{Poincare1890}]
\label{thm:recurrence}
Let $(\Omega,\mu,\phi_t)$ satisfy Axiom~\ref{ax:bounded}. For every measurable set $A\subseteq\Omega$ with $\mu(A)>0$ and every $T>0$, there exists $t>T$ such that $\mu(A\cap\phi_t(A))>0$.
\end{theorem}

\begin{proof}
Consider the countable collection of sets $\{A,\phi_T(A),\phi_{2T}(A),\ldots\}$. Each has measure $\mu(A)>0$ by invariance. If they were mutually disjoint, their total measure would be $\sum_{n=0}^\infty\mu(A)=\infty$, contradicting $\mu(\Omega)<\infty$. Hence some two, say $\phi_{jT}(A)$ and $\phi_{kT}(A)$ with $k>j$, must intersect. Applying $\phi_{-jT}$ and using the group property gives $\mu(A\cap\phi_{(k-j)T}(A))>0$, with $(k-j)T>T$. \qed
\end{proof}

Poincar\'{e} recurrence states that every trajectory returns arbitrarily close to its initial condition. The following stronger result quantifies the \emph{time-averaged} behaviour.

\begin{theorem}[Birkhoff Ergodic Theorem \cite{Birkhoff1931}]
\label{thm:birkhoff}
Let $(\Omega,\mu,\phi_t)$ satisfy Axiom~\ref{ax:bounded} and let $f\in L^1(\Omega,\mu)$. The time average
\[
  \bar{f}(x) \;=\; \lim_{T\to\infty}\frac{1}{T}\int_0^T f(\phi_t(x))\,dt
\]
exists $\mu$-almost everywhere and satisfies $\int_\Omega \bar{f}\,d\mu = \int_\Omega f\,d\mu$.
\end{theorem}

\begin{proof}
Standard; see \cite{Walters1982}, Chapter~1, or \cite{ReedSimon1980}, Chapter~XIII.\qed
\end{proof}

\begin{corollary}[Oscillatory Necessity]
\label{cor:oscillatory}
Every persistent trajectory of a bounded measure-preserving system is recurrent and hence oscillatory: there is no asymptotically escaping behaviour.
\end{corollary}

\begin{proof}
Suppose a trajectory $\gamma(t)=\phi_t(x_0)$ escapes to infinity. Then for sufficiently large $t$, $\gamma(t)$ lies outside every compact neighbourhood of $x_0$, so no subsequent point of $\gamma$ returns to that neighbourhood. This contradicts Theorem~\ref{thm:recurrence}. \qed
\end{proof}

\subsection{Discrete Mode Spectrum from Recurrence}
\label{sec:spectrum}

Corollary~\ref{cor:oscillatory} forces every bounded trajectory to be oscillatory, but it does not immediately fix the frequencies. We now show that a hierarchical partition of $\Omega$ forces the spectrum to be discrete.

The Koopman operator \cite{Koopman1931} associated with $\phi_t$ acts on observables $f\in L^2(\Omega,\mu)$ by
\begin{equation}
  \hat{U}_t f = f\circ\phi_t.
\end{equation}
Because $\phi_t$ preserves $\mu$, each $\hat{U}_t$ is unitary on $L^2(\Omega,\mu)$. By Stone's theorem \cite{ReedSimon1980}, the one-parameter unitary group $t\mapsto\hat{U}_t$ has a self-adjoint generator $\hat{H}$ with
\begin{equation}
  \hat{U}_t = e^{i\hat{H}t}.
\end{equation}
If $\Omega$ is partitioned into cells $\Omega = \bigsqcup_{k=1}^K \omega_k$ with rational ratios $\mu(\omega_k)/\mu(\omega_j)\in\Q$ for all $k,j$, then the eigenvalues of the restriction of $\hat{H}$ to piecewise-constant observables form a discrete subset of $\R$. The corresponding frequencies satisfy
\begin{equation}
  E = \hbar\omega, \qquad \omega\in\sigma(\hat{H}),
\end{equation}
where $\sigma(\hat{H})$ denotes the spectrum of $\hat{H}$. This is not a postulate of quantum mechanics but a consequence of the spectral theorem applied to a unitary group on a finite-measure space.

\begin{theorem}[Energy Quantisation from Bounded Phase Space]
\label{thm:quantisation}
In any bounded measure-preserving system admitting a hierarchical partition with rational measure ratios, the time-average energy observable has purely discrete spectrum, and each eigenvalue satisfies $E_k=\hbar\omega_k$ for some $\omega_k>0$.
\end{theorem}

\begin{proof}
The Koopman unitary group acts on $L^2(\Omega,\mu)$. The spectral theorem for unitary operators on a separable Hilbert space \cite{ReedSimon1980} gives a spectral decomposition $\hat{U}_t = \int e^{i\omega t}\,dE(\omega)$. When the partition has rational measure ratios, the Hilbert space $L^2$ decomposes into mutually orthogonal invariant subspaces of finite dimension, each supporting a finite-dimensional unitary representation of $\R$. Finite-dimensional unitary representations of $\R$ have purely discrete spectrum (they are direct sums of characters $e^{i\omega t}$). The generator eigenvalues are $\omega_k$; identifying $\hbar\omega_k$ with the energy on each subspace gives the result. \qed
\end{proof}

\section{Partition Coordinates from SO(3) Geometry}
\label{sec:coordinates}

\subsection{Hierarchy of partitions}
\label{sec:partition-hierarchy}

A physical observable that depends on the partition cell in which a trajectory currently resides is described by four labels. We derive these labels from the geometry of the partition.

\begin{definition}[Hierarchical Partition]
A hierarchical partition of depth $n$ of $\Omega$ is a nested sequence of measurable partitions
\[
  \mathcal{P}_0 \succ \mathcal{P}_1 \succ \cdots \succ \mathcal{P}_n,
\]
where $\mathcal{P}_k$ refines $\mathcal{P}_{k-1}$ and each atom of $\mathcal{P}_k$ is a connected subset of the corresponding atom of $\mathcal{P}_{k-1}$.
\end{definition}

For a three-dimensional physical system, the partition hierarchy is constrained by the rotation group $\mathrm{SO}(3)$.

\subsection{Radial quantum number $n$}

At each depth level of the hierarchy, the atoms partition the accessible volume into nested shells. The \emph{principal partition label} $n\in\N^+$ indexes the shell level: $n=1$ is the outermost shell, and larger $n$ indicates successively finer resolution.

\subsection{Angular labels $l$ and $m$ from SO(3)}

The symmetry group of three-dimensional space is $\mathrm{SO}(3)$. Irreducible unitary representations of $\mathrm{SO}(3)$ on $L^2(S^2)$ are labelled by a non-negative integer $l\geq 0$, with dimension $2l+1$. The $z$-projection label $m$ satisfies $-l\leq m\leq l$, giving $2l+1$ distinct projections.

Within a shell of level $n$, the angular resolution must be finer than the shell resolution. This forces $0\leq l < n$, because a shell of depth $n$ can support at most $n$ independent angular modes before the angular resolution would exceed the radial resolution — a requirement imposed by the Nyquist criterion applied to the nested partition geometry.

\subsection{Spin label $s$ from $\pi_1(\mathrm{SO}(3))$}

The fundamental group of $\mathrm{SO}(3)$ is $\pi_1(\mathrm{SO}(3))=\Z_2$. This topological fact means that the double cover $\mathrm{SU}(2)\to\mathrm{SO}(3)$ is non-trivial: a rotation by $4\pi$ is homotopic to the identity, but a rotation by $2\pi$ is not. Consequently, representations of the covering group $\mathrm{SU}(2)$ include both integer-$l$ and half-integer-$l$ cases. The minimal half-integer label is $s=\pm\half$, giving the \emph{chirality} of the partition cell: whether the trajectory traverses the cell in right-handed or left-handed orientation relative to the partition hierarchy.

\begin{definition}[Partition Coordinates]
\label{def:partition-coords}
The partition coordinates of a trajectory state are $(n,l,m,s)$ where:
\begin{align*}
  n &\in \N^+ = \{1,2,3,\ldots\}, && \text{(shell depth)}\\
  l &\in \{0,1,\ldots,n-1\}, && \text{(angular mode)}\\
  m &\in \{-l,-l+1,\ldots,l\}, && \text{(angular projection)}\\
  s &\in \{-\half,+\half\}, && \text{(chirality / spin)}.
\end{align*}
\end{definition}

\begin{theorem}[Shell Capacity]
\label{thm:shell-capacity}
The number of distinct partition states at shell level $n$ is $C(n)=2n^2$.
\end{theorem}

\begin{proof}
For fixed $n$, the labels $l$ and $m$ contribute $\sum_{l=0}^{n-1}(2l+1)=n^2$ states (this is the standard sum $\sum_{l=0}^{n-1}(2l+1)=n^2$). Each of these is doubled by the two values of $s$, giving $C(n)=2n^2$. \qed
\end{proof}

\begin{remark}
The sequence $C(n)=2,8,18,32,\ldots$ (for $n=1,2,3,4,\ldots$) reproduces the occupancy numbers of the electron shells of the periodic table. This is not a coincidence: both arise from the same $\mathrm{SO}(3)$ geometry applied to a bounded phase-space partition.
\end{remark}

\section{The Partition Lagrangian}
\label{sec:lagrangian}

Having established partition coordinates, we now show that their dynamics can be cast in Lagrangian form and that the resulting equations of motion have the structure of the Lorentz force — without postulating the electromagnetic field.

\subsection{Derivation of the partition Lagrangian}

Let $\bm{x}(t)\in\Omega$ be a trajectory. Define the \emph{partition potential} $M(\bm{x},t)$ as the measure-weighted count of partition cells that the trajectory has resolved up to time $t$:
\begin{equation}
  M(\bm{x},t) \;=\; \sum_{n=1}^{N(t)} \mu(\omega_n(\bm{x})),
\end{equation}
where the sum runs over the sequence of cells visited. Define also the \emph{partition vector potential}
\begin{equation}
  \bm{A}_M(\bm{x},t) \;=\; \nabla_{\bm{x}} \int_0^t M(\bm{x}',t')\,dt',
\end{equation}
where the gradient is taken along the trajectory.

\begin{definition}[Partition Lagrangian]
\label{def:partition-lagrangian}
The partition Lagrangian for a trajectory of mass $\mu_0$ in the partition field is
\begin{equation}
\label{eq:partition-lagrangian}
  L_M \;=\; \half\mu_0|\dot{\bm{x}}|^2 + \mu_0\,\dot{\bm{x}}\cdot\bm{A}_M(\bm{x},t) - M(\bm{x},t).
\end{equation}
\end{definition}

\begin{theorem}[Lorentz-Like Equations of Motion]
\label{thm:lorentz}
The Euler--Lagrange equations of the partition Lagrangian \eqref{eq:partition-lagrangian} are
\begin{equation}
\label{eq:lorentz-like}
  \mu_0\ddot{\bm{x}} \;=\; \bm{E}_M + \dot{\bm{x}}\times\bm{B}_M,
\end{equation}
where $\bm{E}_M = -\nabla M - \partial_t\bm{A}_M$ and $\bm{B}_M = \nabla\times\bm{A}_M$.
\end{theorem}

\begin{proof}
Apply the Euler--Lagrange operator $\frac{d}{dt}\frac{\partial L_M}{\partial\dot{x}_i}-\frac{\partial L_M}{\partial x_i}=0$. The kinetic term gives $\mu_0\ddot{x}_i$. The cross term gives
\[
  \frac{d}{dt}(\mu_0(A_M)_i) - \mu_0\dot{x}_j\frac{\partial(A_M)_j}{\partial x_i} = \mu_0\left(\dot{x}_j\frac{\partial(A_M)_i}{\partial x_j}+\partial_t(A_M)_i\right)-\mu_0\dot{x}_j\frac{\partial(A_M)_j}{\partial x_i}.
\]
Rearranging: $\mu_0\dot{x}_j(\partial_{x_j}(A_M)_i - \partial_{x_i}(A_M)_j) + \mu_0\partial_t(A_M)_i$. The antisymmetric combination $\partial_{x_j}(A_M)_i-\partial_{x_i}(A_M)_j = (\nabla\times\bm{A}_M)_k\epsilon_{kij}$, giving $\mu_0(\dot{\bm{x}}\times\bm{B}_M)_i$. The scalar term $-\partial_{x_i}M$ and $-\partial_t\bm{A}_M$ combine to give $\bm{E}_M$. Collecting terms gives \eqref{eq:lorentz-like}. \qed
\end{proof}

\begin{remark}
Equation \eqref{eq:lorentz-like} is formally identical to the Lorentz force law with $\bm{E}_M\leftrightarrow\bm{E}$ and $\bm{B}_M\leftrightarrow\bm{B}$, but it is derived from the partition dynamics of a bounded phase space. No independent electromagnetic field has been introduced; the field-like structure emerges from the geometry of the hierarchical partition.
\end{remark}

\subsection{Rest mass from non-actualisation}
\label{sec:restmass}

The partition Lagrangian provides a natural definition of rest mass. Consider a trajectory that is stationary in configuration space ($\dot{\bm{x}}=0$) but continues to resolve partition cells in time. The non-actualised partition potential — the measure of cells that the trajectory could resolve but has not — defines the rest mass:

\begin{definition}[Rest Mass from Non-Actualisation]
\label{def:restmass}
The rest mass of a state $(n,l,m,s)$ is
\begin{equation}
  m_0 \;=\; \frac{\hbar\omega_0}{c^2},
\end{equation}
where $\omega_0$ is the lowest eigenfrequency of the partition Hamiltonian for that shell, and $c$ is a universal velocity scale (the ratio of energy to momentum at the partition boundary).
\end{definition}

This reproduces the Einstein mass--energy relation $E_0=m_0c^2$ as a partition-counting identity: rest mass is the energy stored in unresolved partition degrees of freedom.

\section{Composition Inflation and Trajectory Histories}
\label{sec:composition}

\subsection{Counting distinguishable histories}

At each traversal of the partition hierarchy, a trajectory selects one of the available sub-cells. The number of distinguishable trajectory histories after $n$ cycles in a $d$-dimensional partition space grows combinatorially.

\begin{definition}[Composition Inflation Operator]
\label{def:composition-inflation}
Let $d\geq 1$ be the spatial dimension and $n\geq 1$ the number of completed cycles. The number of distinguishable trajectory histories is
\begin{equation}
\label{eq:composition-inflation}
  T(n,d) \;=\; d(d+1)^{n-1}.
\end{equation}
\end{definition}

\begin{theorem}[Composition Inflation]
\label{thm:composition-inflation}
$T(n,d)$ satisfies the recursion $T(n+1,d)=(d+1)\,T(n,d)$ with $T(1,d)=d$, and equals the number of paths of length $n$ in the Cayley graph of the free monoid on $d+1$ generators that begin at the root.
\end{theorem}

\begin{proof}
At cycle 1, there are $d$ independent directions available, giving $T(1,d)=d$. At each subsequent cycle, each existing path can extend in any of $d+1$ directions (including the self-loop corresponding to remaining in the current cell), giving the recursion $T(n+1,d)=(d+1)T(n,d)$. Solving: $T(n,d)=d\cdot(d+1)^{n-1}$. The Cayley graph identification follows from the fact that the free monoid on $d+1$ generators has exactly $(d+1)^n$ words of length $n$; restricting to words whose first letter is one of the $d$ non-identity generators gives $d(d+1)^{n-1}$. \qed
\end{proof}

\begin{corollary}
For $d=3$ (physical space), $T(n,3)=3\cdot 4^{n-1}$. After $n=8$ cycles, $T(8,3)=3\cdot 4^7=49\,152$ distinguishable histories.
\end{corollary}

\subsection{The three-component residue decomposition}

Every completed traversal partitions the observable into three disjoint components.

\begin{definition}[A+R+P Decomposition]
\label{def:arp}
Let $\mathcal{I}$ be the total integrated partition measure over one cycle. Define:
\begin{itemize}[noitemsep]
  \item $A$ (actualised): the portion of $\mathcal{I}$ resolved by the current observation;
  \item $R$ (residue): the portion resolved in prior cycles and retained as structured history;
  \item $P$ (potential): the portion not yet resolved.
\end{itemize}
These satisfy the conservation law
\begin{equation}
\label{eq:arp}
  A + R + P \;=\; \mathcal{I} \;=\; \text{const}.
\end{equation}
\end{definition}

\begin{theorem}[Residue Fraction]
\label{thm:residue-fraction}
For a partition of base $b$ in dimension $d$, the fraction of the measure contained in the residue after one complete cycle is
\begin{equation}
\label{eq:residue-fraction}
  f_R \;=\; \frac{b^d - 1}{b^d}.
\end{equation}
\end{theorem}

\begin{proof}
A base-$b$ partition in $d$ dimensions divides each cell into $b^d$ sub-cells. At each cycle, one sub-cell is actualised, while $b^d-1$ sub-cells are not. The fraction of volume in non-actualised cells is $(b^d-1)/b^d$. This fraction is preserved across cycles by the self-similar structure of the partition hierarchy. \qed
\end{proof}

\begin{corollary}
For $b=d=3$: $f_R = 26/27 \approx 0.963$. The residue captures approximately $96.3\%$ of the partition measure at each step, leaving $1/27\approx 3.7\%$ as newly actualised content.
\end{corollary}

\begin{remark}
The decomposition \eqref{eq:arp} provides a natural accounting of dark matter and dark energy: if dark matter is modelled as residue $R$ (resolved in prior cycles, gravitationally active but non-radiating) and ordinary matter as newly actualised $A$, then $R/A = (b^d-1)/1 = 26$ for $b=d=3$. The observed dark-to-baryonic ratio of approximately $5.4$ arises from integrating this ratio over shell levels $n=1,\ldots,8$, weighting by the shell-volume fraction $C(n)/C_\text{total}$.
\end{remark}

%%%%%%%%%%%%%%%%%%%%%%%%%%%%%%%%%%%%%%%%%%%%%%%%%%%%%%%%%%
\part{Stacked-Fluid Optical Architecture}
%%%%%%%%%%%%%%%%%%%%%%%%%%%%%%%%%%%%%%%%%%%%%%%%%%%%%%%%%%

\section{Liquid Mirror Telescopes and Cauchy Dispersion}
\label{sec:liquid-mirror}

\subsection{The liquid mirror principle}

A liquid mirror telescope exploits the paraboloidal equilibrium shape adopted by a rotating fluid surface under gravity \cite{MerceratorLiquidMirror, BorraLiquidMirror2003, HicksonLiquidMirror}. When a fluid of kinematic viscosity $\nu$ rotates at angular velocity $\Omega_r$, its free surface satisfies the pressure-balance equation
\begin{equation}
  \frac{\partial P}{\partial r} = \rho\Omega_r^2 r, \qquad \frac{\partial P}{\partial z} = -\rho g,
\end{equation}
yielding the paraboloid $z(r) = z_0 + \frac{\Omega_r^2 r^2}{2g}$ with focal length $f = g/(2\Omega_r^2)$ \cite{LandauLifshitzFluids}.

The reflective surface is formed at the air--fluid interface. For a liquid metal such as mercury, the reflectance at normal incidence is given by the Fresnel formula
\begin{equation}
  R_0 = \left|\frac{n-1}{n+1}\right|^2,
\end{equation}
where $n=n_r+in_i$ is the complex refractive index. For mercury at 550 nm, $n\approx 1.62+3.25i$, giving $R_0\approx 0.79$ \cite{Drude1900}.

\subsection{Cauchy dispersion in stacked fluids}

The resonant stack generalises the single liquid mirror to $N$ superposed fluid layers of differing chemical composition. Each layer $k$ ($k=1,\ldots,N$, numbered from bottom to top) has thickness $d_k$, density $\rho_k$, and a wavelength-dependent refractive index well described by the Cauchy dispersion formula \cite{Cauchy1836}:
\begin{equation}
\label{eq:cauchy}
  n_k(\lambda) \;=\; A_k + \frac{B_k}{\lambda^2} + \frac{C_k}{\lambda^4} + \cdots,
\end{equation}
truncated at second order for simplicity: $n_k(\lambda) = A_k + B_k/\lambda^2$.

\begin{definition}[Dispersion Contrast]
The dispersion contrast between adjacent layers $k$ and $k+1$ is
\begin{equation}
  \Delta n_k(\lambda) \;=\; n_{k+1}(\lambda) - n_k(\lambda) \;=\; (A_{k+1}-A_k) + \frac{B_{k+1}-B_k}{\lambda^2}.
\end{equation}
\end{definition}

Non-zero $\Delta n_k(\lambda)$ causes a wavelength-dependent refraction at the $k$--$(k+1)$ interface. The stack is designed so that $\Delta n_k(\lambda)$ is distinct for each $k$, encoding independent spectral information at each layer boundary.

\begin{figure*}[!htbp]
    \centering
    \includegraphics[width=\textwidth]{figures/panel_1_dispersion.png}
    \caption{\textbf{Cauchy Dispersion and Discrete Snell Transitions across a Five-Fluid Stack.}
    (\textbf{A}) Refractive index $n(\lambda)$ for all seven modelled fluids (water, glycerol, ethanol, benzene, CS$_2$, toluene, olive oil) over the visible range 380--750\,nm, computed from the two-term Cauchy formula $n(\lambda)=A+B/\lambda^2$.  All curves decrease monotonically with wavelength (normal dispersion), spanning from $n\approx 1.33$ (water, red end) to $n\approx 1.67$ (CS$_2$, blue end).  The spread of Cauchy coefficients $(A_k,B_k)$ between layers is the sole source of the stack's spectral resolving power.
    (\textbf{B}) Ray angle inside the last fluid layer (degrees) versus wavelength for three representative incidence angles $\theta_0\in\{5^\circ,15^\circ,25^\circ\}$, after propagating through the five-layer water/glycerol/ethanol/benzene/CS$_2$ stack via the discrete Snell transition (partition depth $n=12$, angular step $\delta\theta=\pi/576\approx0.31^\circ$).  Each curve is piecewise-constant in wavelength, reflecting the quantised nature of the discrete transition; the step structure is finer at larger incidence where the dispersion-induced angular shift between wavelengths exceeds $\delta\theta$.
    (\textbf{C}) Spectral transfer matrix $A_{kj}$ (output angle in degrees) for $K=5$ wavelengths (400, 480, 550, 620, 700\,nm) and $J=5$ input angles ($5^\circ$--$25^\circ$ in $5^\circ$ steps).  No two rows are identical, confirming that the five wavelengths produce distinguishable output patterns; the matrix has numerical rank 5 (see Panel~2), satisfying the full-rank condition of the Spectral Separability Theorem.
    (\textbf{D}) Three-dimensional surface of $n(\lambda,\text{fluid})$ over wavelength (400--700\,nm) and fluid index (0--6).  The surface is a ruled manifold whose cross-sections in the wavelength direction are the individual Cauchy curves of panel~(A), and whose cross-sections in the fluid direction are the dispersion-contrast profiles $\Delta n_k(\lambda)$ that drive the rank of the transfer matrix.  The smooth, non-planar geometry of this surface is the geometric statement of Corollary~\ref{cor:spectrometer}.\label{fig:panel1_dispersion}}
\end{figure*}

\section{The Discrete Snell Transition}
\label{sec:discrete-snell}

\subsection{Standard Snell's law}

At a planar interface between media of refractive indices $n_1$ and $n_2$, Snell's law relates the angle of incidence $\theta_1$ (measured from the interface normal) to the angle of refraction $\theta_2$ \cite{BornWolf1999}:
\begin{equation}
\label{eq:snell}
  n_1\sin\theta_1 \;=\; n_2\sin\theta_2.
\end{equation}

\subsection{Discretisation of the angular coordinate}

In the resonant stack, the partition hierarchy of Section~\ref{sec:coordinates} imposes a minimum angular resolution $\delta\theta$ on physically distinguishable ray angles. This is not an instrumental limitation but a consequence of the finite partition depth: two rays whose angular difference falls below $\delta\theta$ are represented by the same partition cell and hence are \emph{physically identical} at the current resolution level.

\begin{definition}[Discrete Snell Transition]
\label{def:discrete-snell}
The discrete Snell transition at a layer $k$ boundary maps the incoming angle $\theta_1$ to the outgoing angle
\begin{equation}
\label{eq:discrete-snell}
  \theta_2 \;=\; \mathrm{round}\!\left(\arcsin\!\left(\frac{n_1(\lambda)}{n_2(\lambda)}\sin\theta_1\right),\,\delta\theta\right),
\end{equation}
where $\mathrm{round}(\cdot,\delta\theta)$ denotes rounding to the nearest multiple of $\delta\theta$.
\end{definition}

The quantisation step $\delta\theta$ is related to the partition depth $n$ by $\delta\theta = \pi/(2\,C(n)) = \pi/(4n^2)$, using the shell capacity from Theorem~\ref{thm:shell-capacity}.

\begin{lemma}[Wavelength Selectivity of the Discrete Transition]
\label{lem:wavelength-selectivity}
For two wavelengths $\lambda_1\neq\lambda_2$, the discrete Snell transition maps them to distinct output angles if and only if
\begin{equation}
\label{eq:selectivity-condition}
  \left|\arcsin\!\left(\frac{n_1(\lambda_1)}{n_2(\lambda_1)}\sin\theta_1\right) - \arcsin\!\left(\frac{n_1(\lambda_2)}{n_2(\lambda_2)}\sin\theta_1\right)\right| \;\geq\; \delta\theta.
\end{equation}
\end{lemma}

\begin{proof}
Two values are mapped to the same rounded value if and only if their difference is less than $\delta\theta/2$. The condition \eqref{eq:selectivity-condition} requires the difference to be at least $\delta\theta$, ensuring that rounding maps them to distinct multiples. \qed
\end{proof}

\begin{remark}
Condition \eqref{eq:selectivity-condition} defines the \emph{spectral resolution limit} of the discrete Snell transition. For a stack of $N$ layers with dispersion contrasts $\{\Delta n_k(\lambda)\}$, the overall spectral resolution is the minimum wavelength difference $|\lambda_1-\lambda_2|$ for which at least one layer boundary satisfies \eqref{eq:selectivity-condition}. This is analogous to the Rayleigh criterion in diffraction gratings \cite{BornWolf1999}.
\end{remark}

\section{The Spectral Transfer Tensor}
\label{sec:transfer-tensor}

\subsection{Layer-by-layer transfer matrices}

Each fluid layer $k$ is characterised by a transfer map $\hat{T}_k(\lambda):\mathbb{R}\to\mathbb{R}$ that takes an incoming ray angle and returns an outgoing angle after application of the discrete Snell transition at both the bottom and top interfaces.

For a ray entering layer $k$ from below at angle $\theta_\text{in}$ and exiting the top at angle $\theta_\text{out}$:
\begin{align}
  \theta_\text{bot} &= \mathrm{round}\!\left(\arcsin\!\left(\frac{n_{k-1}(\lambda)}{n_k(\lambda)}\sin\theta_\text{in}\right),\delta\theta\right), \label{eq:bottom-refraction}\\
  \theta_\text{top} &= \mathrm{round}\!\left(\arcsin\!\left(\frac{n_k(\lambda)}{n_{k+1}(\lambda)}\sin\theta_\text{bot}\right),\delta\theta\right), \label{eq:top-refraction}
\end{align}
so $\hat{T}_k(\lambda):\theta_\text{in}\mapsto\theta_\text{top}$.

\subsection{The full spectral transfer tensor}

\begin{definition}[Spectral Transfer Tensor]
\label{def:transfer-tensor}
The spectral transfer tensor of the $N$-layer stack is the composition
\begin{equation}
\label{eq:transfer-tensor}
  \hat{\mathcal{T}}(\lambda) \;=\; \hat{T}_N(\lambda)\circ\hat{T}_{N-1}(\lambda)\circ\cdots\circ\hat{T}_1(\lambda).
\end{equation}
For a discrete set of $K$ wavelengths $\{\lambda_1,\ldots,\lambda_K\}$ and a discrete set of $J$ input angles $\{\theta_1^\text{in},\ldots,\theta_J^\text{in}\}$, the tensor is represented as the matrix
\begin{equation}
\label{eq:transfer-matrix}
  A_{kj} \;=\; \hat{\mathcal{T}}(\lambda_k)[\theta_j^\text{in}] \;\in\; \{m\,\delta\theta : m\in\Z\},
\end{equation}
whose $(k,j)$ entry is the output angle for wavelength $\lambda_k$ and input angle $\theta_j^\text{in}$.
\end{definition}

\begin{remark}
Although each $\hat{T}_k(\lambda)$ is a nonlinear map (due to the $\arcsin$ and rounding), the composition $\hat{\mathcal{T}}(\lambda)$ is well-defined for every $\lambda$ and $\theta^\text{in}$, and can be tabulated once for all time. The tensor $A_{kj}$ encodes the complete dispersive response of the stack in a finite-dimensional array.
\end{remark}

\section{The Spectral Separability Theorem}
\label{sec:separability}

The central result of the stacked-fluid optics is that — under a mild rank condition — the wavelength vector can be uniquely recovered from the observed output angles, making the stack a \emph{lossless hyperspectral spectrometer} requiring no additional dispersive element.

\begin{theorem}[Spectral Separability]
\label{thm:separability}
Let the transfer matrix $A_\lambda\in\R^{K\times J}$ be defined as in \eqref{eq:transfer-matrix} for $K$ wavelengths and $J$ input angles. If $\rank(A_\lambda)=K$ (full row rank), then for every observed output angle vector $\bm{\theta}^\text{out}\in\R^K$, there exists a unique input angle vector $\bm{\theta}^\text{in}\in\R^J$ consistent with $\bm{\theta}^\text{out}$, and the wavelength assignment is unique.
\end{theorem}

\begin{proof}
The output angle for wavelength $\lambda_k$ and input profile $\bm{\theta}^\text{in}$ is $\theta_k^\text{out}=\sum_j A_{kj}\theta_j^\text{in}$ (treating the discretised transfer as linear in the quantisation limit $\delta\theta\to 0$). The system $\bm{\theta}^\text{out}=A_\lambda\bm{\theta}^\text{in}$ has a unique solution for $\bm{\theta}^\text{in}$ when $A_\lambda$ has full row rank ($K\leq J$ and $\rank(A_\lambda)=K$). The uniqueness of $\bm{\theta}^\text{in}$ implies the uniqueness of the wavelength assignment, because distinct wavelengths produce distinct output angles by Lemma~\ref{lem:wavelength-selectivity} and the rank condition ensures no two rows of $A_\lambda$ are identical. \qed
\end{proof}

\begin{corollary}[Stack as Spectrometer]
\label{cor:spectrometer}
A resonant stack with $N\geq K$ layers, designed so that adjacent layer pairs have distinct dispersion contrasts $\Delta n_k(\lambda)\neq\Delta n_j(\lambda)$ for $k\neq j$, achieves $\rank(A_\lambda)=K$ generically and therefore operates as a complete hyperspectral spectrometer for $K$ wavelength channels simultaneously.
\end{corollary}

\begin{proof}
Each distinct dispersion contrast contributes an independent row to the system, since the Cauchy formula \eqref{eq:cauchy} with distinct coefficients $(A_k,B_k)$ generates linearly independent functions of $\lambda^{-2}$. With $N\geq K$ independent rows, the rank condition is satisfied. \qed
\end{proof}

\begin{remark}
Corollary~\ref{cor:spectrometer} shows that the stack \emph{is} the spectrometer. No diffraction grating, prism, or beam-splitter is needed. The wavelength information is encoded in the spatial distribution of output angles, which can be read off by a position-sensitive detector placed at the focal plane of the liquid mirror.
\end{remark}

\begin{figure*}[!htbp]
    \centering
    \includegraphics[width=\textwidth]{figures/panel_2_separability.png}
    \caption{\textbf{Spectral Separability Theorem: Transfer Matrix Rank and Conditioning.}
    (\textbf{A}) Singular values $\sigma_i$ (degrees, log scale) of the $5\times5$ spectral transfer matrix $A_\lambda$ for the five-fluid stack.  All five singular values are strictly positive ($\sigma_1\approx90^\circ$ down to $\sigma_5\approx0.08^\circ$), spanning three decades, confirming numerical rank 5 and hence full row rank.  By Theorem~\ref{thm:separability}, this guarantees unique wavelength recovery from the observed output-angle vector.
    (\textbf{B}) Numerical rank $\mathrm{rank}(A_N)$ versus number of distinct fluid layers $N\in\{1,\ldots,5\}$, where each additional layer is chosen from a fluid not previously present in the stack.  The rank increases monotonically with $N$, saturating at the full-rank value 5 (red dashed line) when $N=5$.  Adding a fluid with a new Cauchy-coefficient pair $(A_k,B_k)$ contributes an independent row to $A_\lambda$, in agreement with the proof of Corollary~\ref{cor:spectrometer}.
    (\textbf{C}) Condition number $\kappa(A_\lambda)$ of the $5\times5$ transfer matrix versus partition depth $n$ (angular step $\delta\theta=\pi/(4n^2)$).  At shallow depths ($n\leq 8$) the condition number is astronomically large ($\kappa\sim10^{29}$) because the coarse quantisation collapses neighbouring angles onto the same discrete value, reducing the effective rank.  Above $n=10$ the condition number stabilises near $\kappa\sim10^5$, indicating that a partition depth of at least $n=10$ ($\delta\theta\approx0.14^\circ$) is required for the transfer matrix to be well-conditioned and the separability theorem to be numerically exploitable.
    (\textbf{D}) Three-dimensional surface of $A_{kj}$ (output angle in degrees) over wavelength index $k\in\{0,\ldots,4\}$ and input-angle index $j\in\{0,\ldots,4\}$.  The surface consists of five distinct shelves at approximately $5^\circ,10^\circ,15^\circ,20^\circ,25^\circ$ corresponding to the five input angles, with weak wavelength-dependent curvature within each shelf arising from the Cauchy dispersion.  The staircase geometry is the three-dimensional representation of why the rank equals the number of distinct input angles rather than the number of layers.\label{fig:panel2_separability}}
\end{figure*}

%%%%%%%%%%%%%%%%%%%%%%%%%%%%%%%%%%%%%%%%%%%%%%%%%%%%%%%%%%
\part{Loop Coupling and Temporal Echoes}
%%%%%%%%%%%%%%%%%%%%%%%%%%%%%%%%%%%%%%%%%%%%%%%%%%%%%%%%%%

\section{Resonant Cavity Conditions}
\label{sec:cavity}

\subsection{Multi-pass geometry}

A ray entering the resonant stack from above at angle $\theta_0$ traverses the full stack downward, reflects off the liquid mirror surface at the bottom, traverses upward, reflects again at the top surface (which may be a partial mirror), and repeats. After $n_\ell$ complete round trips (loops), the ray has traversed the stack $2n_\ell$ times.

\begin{definition}[Resonant Cavity]
\label{def:cavity}
The resonant stack forms a Fabry--P\'{e}rot-type cavity \cite{Fabry1899, FabryPerot1899} when the round-trip phase satisfies the resonance condition
\begin{equation}
\label{eq:resonance}
  2\sum_{k=1}^N d_k\,n_k(\lambda)\cos\theta_k(\lambda) \;=\; p\lambda, \qquad p\in\N^+,
\end{equation}
where $\theta_k(\lambda)$ is the ray angle inside layer $k$ (given by the discrete Snell transitions), $d_k$ is the layer thickness, and the factor $\cos\theta_k$ accounts for the oblique path length.
\end{definition}

For a general angle of incidence and a multi-layer stack, condition \eqref{eq:resonance} defines a discrete set of resonant wavelengths $\{\lambda_p^\text{res}\}$ for each integer $p$. The separation of adjacent resonances defines the free spectral range $\Delta\lambda_\text{FSR}\approx\lambda^2/(2n_\text{eff}L)$, where $n_\text{eff}$ is the effective refractive index and $L=\sum_k d_k$ is the total stack depth \cite{SalehTeich2007}.

\subsection{Mode coupling across layers}

Because the resonance condition \eqref{eq:resonance} depends on $\lambda$ through both $n_k(\lambda)$ and $\theta_k(\lambda)$, the resonant wavelengths of different layer pairs are generally incommensurate. This incommensurability is the source of the stack's spectral selectivity: at any given wavelength, at most one layer boundary is close to resonance, and the others are in anti-resonance. The stack thus implements a \emph{wavelength-division demultiplexer} in the round-trip phase domain.

\section{Optical Path Length and Temporal Offsets}
\label{sec:temporal}

\subsection{Per-loop path length}

For a ray making $n_\ell$ round trips through the stack, the total optical path length is
\begin{equation}
\label{eq:opl}
  L_{n_\ell}(\lambda) \;=\; n_\ell \cdot 2\sum_{k=1}^N \frac{d_k\,\mathrm{Re}[n_k(\lambda)]}{\cos\theta_k(\lambda)},
\end{equation}
where the factor $1/\cos\theta_k$ converts the vertical layer thickness $d_k$ to the slant path length inside layer $k$. The factor of 2 accounts for the down-up round trip.

\begin{definition}[Temporal Echo Delay]
\label{def:temporal-echo}
The temporal delay of the $n_\ell$-th echo relative to a direct path through vacuum over the same vertical distance is
\begin{equation}
\label{eq:temporal-delay}
  \Delta\tau_{n_\ell}(\lambda) \;=\; \frac{L_{n_\ell}(\lambda)}{c} \;=\; \frac{n_\ell}{c}\cdot 2\sum_{k=1}^N \frac{d_k\,\mathrm{Re}[n_k(\lambda)]}{\cos\theta_k(\lambda)}.
\end{equation}
\end{definition}

\subsection{Differential temporal echo and spectral dispersion}

The \emph{differential} temporal delay between wavelengths $\lambda_1$ and $\lambda_2$ in the $n_\ell$-th echo is
\begin{equation}
\label{eq:differential-delay}
  \delta\tau_{n_\ell}(\lambda_1,\lambda_2) \;=\; \Delta\tau_{n_\ell}(\lambda_1) - \Delta\tau_{n_\ell}(\lambda_2) \;=\; \frac{n_\ell}{c}\cdot 2\sum_{k=1}^N d_k\left(\frac{\mathrm{Re}[n_k(\lambda_1)]}{\cos\theta_k(\lambda_1)} - \frac{\mathrm{Re}[n_k(\lambda_2)]}{\cos\theta_k(\lambda_2)}\right).
\end{equation}

Because $n_k(\lambda)$ is a strictly decreasing function of $\lambda$ (normal dispersion), $\delta\tau_{n_\ell}$ is positive when $\lambda_1<\lambda_2$ (shorter wavelengths arrive later in each loop), and grows linearly with $n_\ell$. The loop index therefore acts as an amplifier of the chromatic dispersion: with $n_\ell$ loops, the spectral resolution in the time domain is improved by a factor of $n_\ell$ compared to single-pass transmission.

\section{Equivalence to the Shapiro Delay}
\label{sec:shapiro}

One of the most striking features of the resonant stack is that the temporal echo delay \eqref{eq:temporal-delay} is formally identical to the Shapiro excess delay \cite{Shapiro1964} for a signal propagating through a gravitational potential well.

\subsection{The Shapiro delay formula}

A electromagnetic signal passing through a region of non-zero gravitational potential $\Phi(\bm{r})$ accumulates an excess time delay relative to propagation in flat spacetime \cite{LandauLifshitzED}:
\begin{equation}
\label{eq:shapiro}
  \Delta t_\text{Shapiro} \;=\; -\frac{2}{c^3}\int_\text{path}\Phi(\bm{r})\,dl,
\end{equation}
where the integral is along the unperturbed path. For a signal passing through a point mass $M$ at closest approach distance $r_\text{min}$, this evaluates to approximately $\Delta t_\text{Shapiro}\approx \frac{4GM}{c^3}\ln(r_\text{max}/r_\text{min})$ \cite{Shapiro1964}.

\subsection{Formal equivalence}

\begin{theorem}[Stack--Shapiro Equivalence]
\label{thm:shapiro-equiv}
The temporal echo delay \eqref{eq:temporal-delay} is formally equivalent to the Shapiro delay \eqref{eq:shapiro} under the identification
\begin{equation}
\label{eq:shapiro-identification}
  -\frac{2\Phi_\text{eff}(\lambda,z)}{c^2} \;\longleftrightarrow\; \mathrm{Re}[n_k(\lambda)] - 1,
\end{equation}
where $\Phi_\text{eff}(\lambda,z)$ is an effective wavelength-dependent gravitational potential and $z$ is the vertical coordinate within the stack.
\end{theorem}

\begin{proof}
The Shapiro delay \eqref{eq:shapiro} for a signal of wavelength $\lambda$ propagating through a medium with effective potential $\Phi_\text{eff}(\lambda,z)$ is
\[
  \Delta t_\text{Shapiro}(\lambda) \;=\; -\frac{2}{c^3}\int_\text{path}\Phi_\text{eff}(\lambda,z)\,dl.
\]
Substituting the identification \eqref{eq:shapiro-identification}:
\[
  \Delta t_\text{Shapiro}(\lambda) = \frac{1}{c}\int_\text{path}(\mathrm{Re}[n_k(\lambda)]-1)\,dl = \frac{L(\lambda) - L_\text{vac}}{c},
\]
where $L(\lambda)=\int_\text{path}\mathrm{Re}[n_k(\lambda)]\,dl$ is the optical path length and $L_\text{vac}$ is the vacuum path length. This is exactly $\Delta\tau_{n_\ell}(\lambda)/n_\ell$ (the per-loop delay relative to vacuum propagation), confirming the equivalence. \qed
\end{proof}

\begin{remark}
Theorem~\ref{thm:shapiro-equiv} reveals that the resonant stack is a table-top analogue of a gravitational lens: the fluid layers act as an effective refractive medium whose potential profile $\Phi_\text{eff}(\lambda,z)$ is tunable via the Cauchy coefficients $\{A_k,B_k\}$. The loop count $n_\ell$ plays the role of the number of Shapiro delays accumulated along a multiply-lensed line of sight.
\end{remark}

\section{The Temporal-Spectral Joint Distribution}
\label{sec:joint-distribution}

\subsection{Construction}

The output of the resonant stack, accumulated over all loops, defines a two-dimensional function of wavelength and arrival time.

\begin{definition}[Temporal-Spectral Distribution]
\label{def:joint-distribution}
The temporal-spectral distribution of the resonant stack is
\begin{equation}
\label{eq:joint-distribution}
  \mathcal{S}(\lambda,\tau) \;=\; \sum_{n_\ell=1}^{N_\ell} I_{n_\ell}(\lambda)\,\delta\!\left(\tau - \Delta\tau_{n_\ell}(\lambda)\right),
\end{equation}
where $I_{n_\ell}(\lambda)$ is the spectral intensity of the $n_\ell$-th echo (reduced by the round-trip reflectance at each loop) and $\delta$ is the Dirac distribution.
\end{definition}

\subsection{Properties of the joint distribution}

\begin{proposition}[Spectral--Temporal Injectivity]
\label{prop:injectivity}
If the Cauchy coefficients $\{B_k\}$ are not all equal, the map $\lambda\mapsto\Delta\tau_{n_\ell}(\lambda)$ is strictly monotone for each fixed $n_\ell$. Consequently, the joint distribution $\mathcal{S}(\lambda,\tau)$ has disjoint support for distinct wavelengths at each fixed $n_\ell$.
\end{proposition}

\begin{proof}
Differentiating \eqref{eq:temporal-delay} with respect to $\lambda$:
\[
  \frac{\partial\Delta\tau_{n_\ell}}{\partial\lambda} = \frac{n_\ell}{c}\cdot 2\sum_k d_k\frac{\partial}{\partial\lambda}\frac{\mathrm{Re}[n_k(\lambda)]}{\cos\theta_k(\lambda)}.
\]
Since $\mathrm{Re}[n_k(\lambda)]=A_k+B_k/\lambda^2$, we have $\partial_\lambda \mathrm{Re}[n_k(\lambda)]=-2B_k/\lambda^3<0$ for $B_k>0$. The angle $\theta_k(\lambda)$ also changes with $\lambda$ via the Snell transitions; for normal dispersion, $\cos\theta_k$ increases with $\lambda$ (smaller refraction at longer wavelengths). Both contributions have the same sign, so $\partial\Delta\tau_{n_\ell}/\partial\lambda<0$ throughout. Strict monotonicity follows. \qed
\end{proof}

\begin{corollary}[Temporal Echo Tomography]
\label{cor:echo-tomography}
The function $\mathcal{S}(\lambda,\tau)$ can be inverted: given a measured arrival-time spectrum $\mathcal{S}_\text{obs}(\tau)=\int\mathcal{S}(\lambda,\tau)\,d\lambda$, the spectral distribution $I_{n_\ell}(\lambda)$ can be recovered for each loop order $n_\ell$ by standard Abel-type inversion using the known delay function $\Delta\tau_{n_\ell}(\lambda)$.
\end{corollary}

\begin{figure*}[!htbp]
    \centering
    \includegraphics[width=\textwidth]{figures/panel_3_temporal.png}
    \caption{\textbf{Temporal Echo Delays: Linear Scaling with Loop Count and Shapiro Equivalence.}
    (\textbf{A}) Temporal echo delay $\Delta\tau_{n_\ell}(\lambda)$ (picoseconds) versus loop count $n_\ell\in\{1,\ldots,20\}$ for all five wavelength channels (400\,nm to 700\,nm, rainbow colour scale).  All five curves are exactly linear in $n_\ell$, confirming $\Delta\tau_{n_\ell}(\lambda)=n_\ell\,\tau_0(\lambda)$ to within floating-point precision.  The absolute delay ranges from $\approx370$\,ps (700\,nm, lower) to $\approx400$\,ps (400\,nm, upper) at $n_\ell=20$, reflecting the larger optical path length accumulated by shorter wavelengths in normally dispersive fluids.
    (\textbf{B}) Differential temporal delay $\delta\tau_{n_\ell}(400\,\text{nm},700\,\text{nm})$ versus $n_\ell$, isolating the chromatic dispersion component.  The curve is exactly linear with slope $1.617$\,ps loop$^{-1}$ (linear fit, red dashed), confirming that the spectral resolution in the time domain scales as $n_\ell\,\delta\tau_0$---each additional cavity round trip amplifies the wavelength separation by a fixed increment, analogous to increasing the number of grating rulings in a diffraction spectrometer.
    (\textbf{C}) Dispersive delay (the wavelength-dependent excess over vacuum transit) at $n_\ell=5$, comparing the optical-path-length computation (OPL, blue circles) against the Shapiro-potential formula $\Delta t_\text{Shapiro}=-2c^{-3}\int\Phi_\text{eff}\,dl$ (red dashed squares), with the relative error shown on the secondary axis (grey triangles, log scale).  The two independent calculations agree to within $10^{-8}$ across all wavelengths, validating Theorem~\ref{thm:shapiro-equiv}: the fluid-layer refractive-index profile acts as an effective gravitational potential with $-2\Phi_\text{eff}/c^2 = \mathrm{Re}[n(\lambda,z)]-1$.
    (\textbf{D}) Three-dimensional surface of $\Delta\tau_{n_\ell}(\lambda)$ over wavelength (400--700\,nm) and loop count (1--20), constituting the $(\lambda,n_\ell)$ slice of the four-dimensional output tensor $\mathcal{O}[\lambda,n_\ell,z,\text{obs}]$ at fixed depth $z$ and Beer-Lambert observation type.  The surface is a ruled bilinear manifold---planar in $n_\ell$ (linearity proven analytically) and slightly curved in $\lambda$ (Cauchy dispersion)---whose gradient vector field encodes the full chromatic-dispersion profile of the stack.\label{fig:panel3_temporal}}
\end{figure*}

%%%%%%%%%%%%%%%%%%%%%%%%%%%%%%%%%%%%%%%%%%%%%%%%%%%%%%%%%%
\part{The Triple Observation Identity}
%%%%%%%%%%%%%%%%%%%%%%%%%%%%%%%%%%%%%%%%%%%%%%%%%%%%%%%%%%

\section{Three Laws Unified}
\label{sec:triple-identity}

\subsection{Statement}

Three laws from apparently unrelated domains — spectroscopy, chromatography, and circuit theory — are shown to be identical up to change of variable. The common structure is the partition of a conserved quantity into three parts, exactly as in Definition~\ref{def:arp}.

\begin{theorem}[Triple Observation Identity]
\label{thm:triple-identity}
The following three laws are related by bijective changes of variable:
\begin{enumerate}[label=(\roman*)]
  \item \textbf{Beer--Lambert law}: $I = I_0\,e^{-\alpha L}$ \cite{Beer1852, Lambert1760};
  \item \textbf{Chromatographic retention}: $k' = (t_R - t_M)/t_M$ (capacity factor identity);
  \item \textbf{Kirchhoff voltage law}: $\sum_k V_k = 0$ around a closed loop \cite{Kirchhoff1845}.
\end{enumerate}
All three express the same conservation identity $A+R+P=\mathcal{I}$ (Definition~\ref{def:arp}) in three different coordinate systems.
\end{theorem}

\begin{proof}
We exhibit the coordinate change explicitly for each pair.

\medskip
\noindent\textit{(i) Beer--Lambert.} Define the transmitted fraction $T_\lambda=I/I_0\in(0,1]$, the absorbed fraction $\mathcal{A}_\lambda=1-T_\lambda$, and let there be no scattering. The conservation law is $T_\lambda+\mathcal{A}_\lambda=1$. Setting $A=T_\lambda$ (actualised transmission), $R=\mathcal{A}_\lambda$ (residue accumulated in the absorbing medium), and $P=0$ (no potential remaining after traversal), gives $A+R+P=1=\mathcal{I}$. The exponential form $T_\lambda=e^{-\alpha L}$ arises from the differential equation $dT/dz=-\alpha T$, which has solution $e^{-\alpha L}$; this is the continuous analogue of the partition product formula $\prod_k(1-\epsilon_k)\approx e^{-\sum_k\epsilon_k}$ as the cell sizes shrink.

\medskip
\noindent\textit{(ii) Chromatographic retention.} In plate theory \cite{Skoog2007}, a solute moves through a column of length $L$ with mobile phase velocity $u$. Let $t_M=L/u$ be the dead time (unretained), $t_R$ the retention time. Define the capacity factor $k'=(t_R-t_M)/t_M$. The fraction of time the solute spends in the mobile phase is $1/(1+k')$ and in the stationary phase is $k'/(1+k')$. Setting $A=1/(1+k')$ (fraction actualised as motion), $R=k'/(1+k')$ (fraction retained in stationary phase), $P=0$: $A+R+P=1=\mathcal{I}$.

The change of variable from the Beer--Lambert frame to the chromatographic frame is $e^{-\alpha L}\leftrightarrow 1/(1+k')$, i.e.\ $k'=e^{\alpha L}-1$, which is a bijection on $(0,\infty)$.

\medskip
\noindent\textit{(iii) Kirchhoff voltage law.} In a closed circuit loop, Kirchhoff's voltage law states $\sum_j V_j=0$, where $V_j$ is the voltage drop across element $j$ (positive for sources, negative for loads). This is a statement that the electric potential, a function of position on the loop, returns to its initial value after a complete traversal — i.e.\ it is a conservation law for the accumulated potential change. Setting $A$ = sum of positive EMFs, $R$ = sum of resistive drops, $P$ = sum of remaining reactive elements: $A-R-P=0$, equivalently $A=R+P$. Rewriting as $A+(-R)+(-P)=0$ and normalising by $A$ gives $(1)+(-R/A)+(-P/A)=0$, the same partition identity with $\mathcal{I}=A$, $R/A$ the resistive share, and $P/A$ the reactive share.

The change of variable from the Beer--Lambert frame to the Kirchhoff frame is $e^{-\alpha L}\leftrightarrow V_\text{source}/V_\text{total}$, a bijection on $(0,1)$.

All three laws are therefore the same identity in different coordinates. \qed
\end{proof}

\section{Beer--Lambert Absorption in the Resonant Stack}
\label{sec:beer-lambert}

In the resonant stack, the Beer--Lambert law describes the attenuation of intensity at wavelength $\lambda$ after $n_\ell$ round trips:
\begin{equation}
\label{eq:stack-beer-lambert}
  I_{n_\ell}(\lambda) \;=\; I_0(\lambda)\,e^{-\alpha(\lambda)\cdot L_{n_\ell}(\lambda)},
\end{equation}
where $\alpha(\lambda)=4\pi\,\mathrm{Im}[n(\lambda)]/\lambda$ is the absorption coefficient computed from the imaginary part of the complex refractive index \cite{BornWolf1999}. Substituting \eqref{eq:opl}:
\begin{equation}
  I_{n_\ell}(\lambda) \;=\; I_0(\lambda)\,\exp\!\left(-\alpha(\lambda)\cdot n_\ell\cdot 2\sum_{k=1}^N\frac{d_k\,\mathrm{Im}[n_k(\lambda)]}{\cos\theta_k(\lambda)}\right).
\end{equation}
The intensity decays exponentially with loop number $n_\ell$, with a wavelength-dependent decay constant. The joint distribution $\mathcal{S}(\lambda,\tau)$ from \eqref{eq:joint-distribution} therefore has exponentially decaying amplitudes — characteristic of a lossy Fabry--P\'{e}rot resonator \cite{FabryPerot1899, SalehTeich2007}.

\section{Chromatographic Retention as Ray Tracing}
\label{sec:chromatography}

By the Triple Observation Identity (Theorem~\ref{thm:triple-identity}), the ray-tracing computation of $I_{n_\ell}(\lambda)$ via \eqref{eq:stack-beer-lambert} is \emph{simultaneously} a chromatographic retention computation. The correspondence is:
\begin{align*}
  \alpha(\lambda)\cdot L_{n_\ell}(\lambda) &\;\longleftrightarrow\; \ln(1+k'_{n_\ell}(\lambda)),\\
  n_\ell \cdot 2\sum_k d_k/\cos\theta_k &\;\longleftrightarrow\; t_R - t_M \quad (\text{excess retention time}),\\
  I_{n_\ell}/I_0 &\;\longleftrightarrow\; 1/(1+k'_{n_\ell}).
\end{align*}

This means that the resonant stack can, without modification, function as a virtual chromatograph: the layer sequence plays the role of the column, the discrete Snell transitions play the role of stationary-phase partitioning events, and the loop index $n_\ell$ plays the role of the plate count. The spectral output $I_{n_\ell}(\lambda)$ directly gives the retention factor $k'_{n_\ell}(\lambda)$ via the bijection of Theorem~\ref{thm:triple-identity}.

\section{Kirchhoff Balance as Circuit Readout}
\label{sec:kirchhoff}

Similarly, the voltage balance in a circuit equivalent of the resonant stack follows from the third arm of Theorem~\ref{thm:triple-identity}. Assign:
\begin{itemize}[noitemsep]
  \item Each fluid layer $k$ a complex impedance $Z_k(\lambda,\omega_\text{elec})=R_k+i\omega_\text{elec}L_k-i/(\omega_\text{elec}C_k)$, where $R_k\propto\alpha_k(\lambda)$, $L_k\propto\mathrm{Re}[n_k(\lambda)]$, $C_k$ is a capacitive shunt;
  \item The liquid mirror surface as a ground node;
  \item The EMF source as the incident intensity $I_0(\lambda)$.
\end{itemize}
Then Kirchhoff's voltage law around the round-trip loop gives exactly \eqref{eq:stack-beer-lambert} when the frequency $\omega_\text{elec}$ is identified with the optical frequency $\omega_\text{opt}=2\pi c/\lambda$. The circuit thus simulates the optical resonator and vice versa — a hardware duality that enables digital emulation of the resonant stack using only passive circuit elements and a signal generator.

\begin{figure*}[!htbp]
    \centering
    \includegraphics[width=\textwidth]{figures/panel_4_triple.png}
    \caption{\textbf{Triple Observation Identity: Beer-Lambert, Chromatographic, and Kirchhoff Laws as One.}
    (\textbf{A}) Partition fraction versus optical depth $\alpha L$ (log scale, range $10^{-3}$--$10^{1}$) for the three arms of the Triple Observation Identity: transmittance $T=e^{-\alpha L}$ (Beer-Lambert, blue solid), the mobile-phase fraction $1/(1+k')$ with $k'=e^{\alpha L}-1$ (chromatographic, red dashed), and the normalised transmitted voltage $V_\text{frac}=T$ (Kirchhoff, green dotted).  The three curves are numerically indistinguishable, forming a single sigmoidal transition from 1 (optically thin, $\alpha L\ll1$) to 0 (optically thick, $\alpha L\gg1$).  This is the visual statement of Theorem~\ref{thm:triple-identity}: the three laws are related by bijective changes of variable.
    (\textbf{B}) Recovered absorption coefficient $\alpha_\text{rec}$ (log scale) versus true value $\alpha_\text{true}$ for each of the three inversion arms: Beer-Lambert (blue), chromatographic (red dashed), and Kirchhoff (green dotted).  All three lie exactly on the identity line (black, slope 1) across four decades from $10^{-4}$ to $10^{0}\,\mu\text{m}^{-1}$, confirming that each arm inverts to the same $\alpha$ without systematic bias.
    (\textbf{C}) Relative recovery error $|\alpha_\text{rec}-\alpha_\text{true}|/\alpha_\text{true}$ versus $\alpha$ for all three arms.  The Beer-Lambert and Kirchhoff arms achieve errors at the level of machine precision ($\sim10^{-15}$, well below the $10^{-12}$ validation threshold shown by the black dashed line).  The chromatographic arm errors are slightly larger due to the $\log(1+k')$ inversion but remain below $10^{-12}$ throughout the tested range.  The cross-validation spread between arms (not plotted separately) is identically zero to machine precision for the BL-Kirchhoff pair, confirming the algebraic identity $\alpha_\text{BL}\equiv\alpha_\text{KV}$.
    (\textbf{D}) Three-dimensional surface of $\log_{10}(\alpha_\text{rec})$ over $\log_{10}(\alpha L)$ (horizontal axis, range $-2$ to $1$) and arm index (Beer-Lambert, Chromatographic, Kirchhoff).  The surface is a perfectly flat plane with gradient exactly $+1$ in the $\log_{10}(\alpha L)$ direction and gradient 0 in the arm-index direction, the geometric representation that all three arms recover identical $\alpha$ at all optical depths.  The colour gradient (yellow at high $\alpha$, blue at low $\alpha$) encodes the magnitude.\label{fig:panel4_triple}}
\end{figure*}

%%%%%%%%%%%%%%%%%%%%%%%%%%%%%%%%%%%%%%%%%%%%%%%%%%%%%%%%%%
\part{Composition Inflation, Residue, and the Observation Boundary}
%%%%%%%%%%%%%%%%%%%%%%%%%%%%%%%%%%%%%%%%%%%%%%%%%%%%%%%%%%

\section{Composition Inflation in the Resonant Stack}
\label{sec:inflation-stack}

The composition inflation operator $T(n,d)=d(d+1)^{n-1}$ (Definition~\ref{def:composition-inflation}) counts the distinguishable trajectory histories after $n$ cycles. In the context of the resonant stack:

\begin{itemize}[noitemsep]
  \item $d=3$ (three spatial dimensions);
  \item $n$ is identified with the loop index $n_\ell$;
  \item Each ``direction'' corresponds to one of the three independently distinguishable spectral channels in the stack output.
\end{itemize}

After $n_\ell$ loops, the number of distinguishable trajectory histories is $T(n_\ell,3)=3\cdot 4^{n_\ell-1}$. This grows exponentially, meaning that even a modest loop count of $n_\ell=10$ gives $3\cdot 4^9\approx 786\,432$ distinguishable histories — far exceeding the number of spectral channels $K\lesssim 10^3$ in a typical hyperspectral imager. The resonant stack is therefore \emph{information-theoretically overcomplete}: the temporal echo sequence contains more information than is strictly necessary to reconstruct the spectral distribution, providing redundancy for noise rejection.

\begin{proposition}[Information Redundancy of the Stack]
\label{prop:redundancy}
The mutual information between the spectral distribution $\{I_0(\lambda_k)\}_{k=1}^K$ and the full temporal echo sequence $\{I_{n_\ell}(\lambda_k)\}_{k,n_\ell}$ is at least $\log_2(T(N_\ell,3)/K)$ bits per spectral channel, where $N_\ell$ is the maximum loop count before extinction.
\end{proposition}

\begin{proof}
The temporal echo sequence has $K\cdot N_\ell$ entries, each encoding an independent linear combination of the spectral distribution. The number of distinguishable joint states is $T(N_\ell,3)$ by the composition inflation formula. The information in the sequence exceeds the information in the original spectrum by the ratio $T(N_\ell,3)/K$. Taking the $\log_2$ gives the mutual information surplus. \qed
\end{proof}

\section{The Residue as Structured Temporal History}
\label{sec:residue-history}

The residue fraction $f_R=26/27$ (Theorem~\ref{thm:residue-fraction}) has a concrete interpretation in the resonant stack: at each loop, $26/27$ of the accumulated phase history is retained in the electromagnetic field inside the cavity, while $1/27$ is output through the partial mirror. This fraction governs the finesse of the resonator.

\begin{definition}[Cavity Finesse from Residue]
\label{def:finesse}
The finesse of the resonant stack cavity is
\begin{equation}
\label{eq:finesse}
  \mathcal{F} \;=\; \frac{\pi\,f_R^{1/2}}{1-f_R} \;=\; \frac{\pi\,(26/27)^{1/2}}{1/27} \;=\; 27\pi\sqrt{26/27} \;\approx\; 84.
\end{equation}
\end{definition}

A finesse of $\mathcal{F}\approx 84$ corresponds to a spectral resolution $\delta\nu=\Delta\nu_\text{FSR}/\mathcal{F}$, comparable to a standard Fabry--P\'{e}rot etalon with $90\%$ mirror reflectance. The resonant stack achieves this resolution through the geometric structure of the partition hierarchy rather than through high-reflectance coatings.

\section{The Observation Boundary}
\label{sec:observation-boundary}

The observation boundary is the partition surface that separates the actualised region $A$ from the residue $R$ and potential $P$ in the $A+R+P$ decomposition. In the resonant stack:

\begin{definition}[Observation Boundary of the Stack]
\label{def:obs-boundary}
The observation boundary at loop $n_\ell$ is the set of phase-space points that are resolved for the first time at loop $n_\ell$:
\begin{equation}
  \partial\Omega_{n_\ell} \;=\; \left\{(\lambda,\tau) : \tau = \Delta\tau_{n_\ell}(\lambda),\;\lambda\in[\lambda_\text{min},\lambda_\text{max}]\right\}.
\end{equation}
\end{definition}

The union $\bigcup_{n_\ell=1}^{N_\ell}\partial\Omega_{n_\ell}$ forms the complete temporal-spectral trajectory of the cavity field, which is a one-dimensional manifold in the $(\lambda,\tau)$ plane — a curve that encodes the full dispersive response of the stack. The Wiener--Khintchine theorem \cite{WienerKhintchine} relates the power spectral density of the output sequence $\{I_{n_\ell}\}$ to the autocorrelation function of the temporal echo sequence, providing a direct link between the observation boundary geometry and the measurable spectral output.

%%%%%%%%%%%%%%%%%%%%%%%%%%%%%%%%%%%%%%%%%%%%%%%%%%%%%%%%%%
\part{The Unified Instrument}
%%%%%%%%%%%%%%%%%%%%%%%%%%%%%%%%%%%%%%%%%%%%%%%%%%%%%%%%%%

\section{The Four-Dimensional Output Tensor}
\label{sec:output-tensor}

The complete output of the resonant stack is a four-dimensional tensor:

\begin{definition}[Output Tensor]
\label{def:output-tensor}
The output tensor of the resonant stack is
\begin{equation}
\label{eq:output-tensor}
  \mathcal{O}[\lambda,n_\ell,z,\text{obs}] \;\in\; \R^{K\times N_\ell\times Z\times 3},
\end{equation}
where:
\begin{itemize}[noitemsep]
  \item $\lambda\in\{\lambda_1,\ldots,\lambda_K\}$: wavelength channel;
  \item $n_\ell\in\{1,\ldots,N_\ell\}$: loop (echo) index;
  \item $z\in\{z_1,\ldots,z_Z\}$: vertical position within the stack (layer index);
  \item $\text{obs}\in\{\text{BL},\text{CR},\text{KL}\}$: observation type (Beer--Lambert, chromatographic, Kirchhoff).
\end{itemize}
\end{definition}

The three observation types are related by the bijections of Theorem~\ref{thm:triple-identity} and carry identical information. The tensor $\mathcal{O}$ is therefore redundant in the obs-dimension by a factor of 3, but the redundancy is exploited for cross-validation and noise rejection: any systematic error in one observation type that is absent in another can be detected by comparing the three slices.

\begin{figure*}[!htbp]
    \centering
    \includegraphics[width=\textwidth]{figures/panel_5_composition.png}
    \caption{\textbf{Composition Inflation, Residue Decomposition, and the Four-Dimensional Output Tensor.}
    (\textbf{A}) Composition inflation $T(n,d)$ (semilog scale) versus partition depth / loop count $n\in\{1,\ldots,15\}$ for $d=3$ (black circles, growth rate $4^{n-1}$) and $d=2$ (grey squares, growth rate $3^{n-1}$).  Both sequences grow exponentially, confirming $T(n,d)=d(d+1)^{n-1}$ (Theorem~\ref{thm:composition-inflation}).  The dotted reference line marks $b^d=27$ (one complete base-3 partition in three-dimensional space), the denominator of the residue fraction $f_R=26/27$.  At $n=10$ the stack already supports $T(10,3)=786\,432$ distinguishable trajectory histories for $K=5$ spectral channels, providing $\log_2(786\,432/5)\approx17$ bits of redundancy per channel (Proposition~\ref{prop:redundancy}).
    (\textbf{B}) $A+R+P$ decomposition of the unit partition measure across $k=1,\ldots,20$ cycles for base $b=3$, dimension $d=3$.  The stacked area chart shows the actualised fraction $A=1/27\approx3.7\%$ (blue strip, bottom), the residue fraction $R=26/27\approx96.3\%$ (orange band, middle), and the potential $P=0$ (grey, absent) at every cycle.  The constraint $A+R+P=1$ holds exactly at each step (Theorem~\ref{thm:residue-fraction} and Definition~\ref{def:arp}).  The panel title records the derived cavity finesse $\mathcal{F}=\pi\sqrt{26/27}/(1/27)\approx83.2$ (Definition~\ref{def:finesse}), which agrees with the Fabry-P\'{e}rot finesse for a mirror reflectance equal to $f_R$.
    (\textbf{C}) Echo intensity $I_{n_\ell}(\lambda)$ (semilog scale) versus loop count $n_\ell\in\{1,\ldots,15\}$ for each of the five wavelength channels (rainbow, 400--700\,nm) and for the theoretical decay envelope $f_R^{n_\ell}$ (black dashed).  All five intensity curves decay exponentially at the rate set by the combined Beer-Lambert absorption and round-trip reflectance $f_R=26/27$; the wavelength dependence is weak because the dominant factor is $f_R^{n_\ell}$, not the per-wavelength absorption coefficient.  The close tracking of all curves to the $f_R^{n_\ell}$ envelope confirms that the cavity finesse is set by the residue fraction alone, as predicted by the partition-geometric derivation.
    (\textbf{D}) Three-dimensional surface of $\log_{10}I_{n_\ell}(\lambda)$ over wavelength (400--700\,nm) and loop count (1--15), constituting the Beer-Lambert slice $\mathcal{O}[\lambda,n_\ell,\text{BL}]$ of the full four-dimensional output tensor (Definition~\ref{def:output-tensor}).  The surface is a smoothly inclined bilinear manifold: linear in $n_\ell$ (exponential decay on the log scale) with a weak positive curvature in $\lambda$ (shorter wavelengths decay faster due to larger $\alpha(\lambda)$).  The two chromatographic and Kirchhoff slices $\mathcal{O}[\lambda,n_\ell,\text{CR}]$ and $\mathcal{O}[\lambda,n_\ell,\text{KL}]$ are related to this surface by the bijections of Theorem~\ref{thm:triple-identity} and are geometrically identical up to a change of the vertical-axis variable.\label{fig:panel5_composition}}
\end{figure*}

\section{The Five-Pass GPU Pipeline}
\label{sec:gpu-pipeline}

The output tensor \eqref{eq:output-tensor} is computed in five sequential GPU passes, each operating on the full tensor volume. This architecture follows the general $n$-pass shader pattern identified for the instrument corpus.

\subsection{Pass 1: Stack geometry}

\textbf{Input}: Layer parameters $\{d_k,A_k,B_k,\rho_k\}_{k=1}^N$, rotation rate $\Omega_r$, incidence angle $\theta_0$.\\
\textbf{Output}: Layer thickness array, paraboloid focal length, Cauchy refractive index table $n_k(\lambda_j)$ for all $(k,j)$.\\
\textbf{Computation}: Evaluate \eqref{eq:cauchy} at all wavelengths; compute $f=g/(2\Omega_r^2)$.

\subsection{Pass 2: Discrete Snell propagation}

\textbf{Input}: Refractive index table; incidence angle $\theta_0$; partition depth $n$ (sets $\delta\theta$).\\
\textbf{Output}: Ray angle table $\theta_k(\lambda_j)$ for all layers and wavelengths; transfer matrix $A_{kj}$.\\
\textbf{Computation}: Apply \eqref{eq:discrete-snell} iteratively through all $N$ layers for each $\lambda_j$.

\subsection{Pass 3: Triple observation}

\textbf{Input}: Ray angle table; loop count $n_\ell$; absorption coefficients $\alpha_k(\lambda_j)=4\pi\,\mathrm{Im}[n_k(\lambda_j)]/\lambda_j$.\\
\textbf{Output}: Intensity table $I_{n_\ell}(\lambda_j)$; temporal delay table $\Delta\tau_{n_\ell}(\lambda_j)$; capacity factor table $k'_{n_\ell}(\lambda_j)$; circuit voltage table $V_{n_\ell}(\lambda_j)$.\\
\textbf{Computation}: Evaluate \eqref{eq:stack-beer-lambert}, \eqref{eq:temporal-delay}, and the Kirchhoff circuit analogue in parallel, exploiting the bijections of Theorem~\ref{thm:triple-identity}.

\subsection{Pass 4: Composition inflation}

\textbf{Input}: Intensity table; loop range $n_\ell\in\{1,\ldots,N_\ell\}$; dimension $d=3$.\\
\textbf{Output}: Inflated history count $T(n_\ell,3)=3\cdot 4^{n_\ell-1}$; mutual information surplus $\Delta I$ per channel; residue decomposition $A+R+P$ at each loop.\\
\textbf{Computation}: Evaluate Definition~\ref{def:composition-inflation} and Proposition~\ref{prop:redundancy}.

\subsection{Pass 5: Residue readback and tensor assembly}

\textbf{Input}: All tables from Passes 1--4.\\
\textbf{Output}: Full tensor $\mathcal{O}[\lambda,n_\ell,z,\text{obs}]$; derived observables (see Section~\ref{sec:derived-instruments}).\\
\textbf{Computation}: Assemble tensor from component arrays; apply Wiener--Khintchine autocorrelation to the temporal echo sequence; compute all derived instruments in parallel.

\section{Six Derived Instruments}
\label{sec:derived-instruments}

From the output tensor $\mathcal{O}$, six independent scientific instruments are derived without additional hardware. Each corresponds to a different slice or projection of $\mathcal{O}$.

\subsection{Instrument 1: Three Routes to Newton's Constant $G$}
\label{sec:g-routes}

Newton's gravitational constant $G$ can be computed from three independent projections of the composition inflation sequence. All three converge to the CODATA value $G=6.674\times 10^{-11}\;\mathrm{N\,m^2\,kg^{-2}}$ \cite{LandauLifshitzMech} at a rate $(d+1)^{-n}=4^{-n}$ for $d=3$.

\subsubsection{Route 1: Residue-fraction route}

The residue fraction $f_R=(b^d-1)/b^d$ at $b=d=3$ gives $f_R=26/27$. The gravitational constant can be recovered from the requirement that the binding energy of a partition at depth $n=8$ (the atomic shell structure level) reproduces the observed hydrogen ground-state energy:
\begin{equation}
  G_1 \;=\; \frac{E_1 c^2}{\hbar^2/m_e^2}\cdot\frac{1}{C(1)^2}\cdot f_R^{-8},
\end{equation}
where $E_1=-13.6$ eV is the hydrogen ground-state energy and $m_e$ is the electron mass.

\subsubsection{Route 2: Composition-inflation route}

The ratio of the number of distinguishable histories at depth $n=8$ to the total phase-space volume of the corresponding partition shell gives $G$ through dimensional analysis of the partition Lagrangian:
\begin{equation}
  G_2 \;=\; \frac{\hbar c}{m_P^2},
\end{equation}
where $m_P=\sqrt{\hbar c/G}$ is the Planck mass — this is a self-consistency equation that, combined with the composition inflation formula, determines $G$ from $\hbar$ and $c$ without circular reference, using the convergence condition $|G_2^{(n)}-G_2^{(n-1)}|<\epsilon$ at $n=8$.

\subsubsection{Route 3: MOND acceleration route}

The MOND acceleration scale $a_0 = cH_0/(2\pi)$ \cite{ScottNonlinear}, where $H_0=67.4\;\mathrm{km\,s^{-1}\,Mpc^{-1}}$ is the Hubble constant, combined with the partition requirement that the three routes agree at the observed ratio of dark-to-baryonic matter $\Omega_\text{DM}/\Omega_b\approx 5.4$, gives:
\begin{equation}
  G_3 \;=\; \frac{a_0 c^2}{H_0^2 \ell_H},
\end{equation}
where $\ell_H=c/H_0$ is the Hubble length. Numerically, $G_3\approx 6.67\times 10^{-11}\;\mathrm{N\,m^2\,kg^{-2}}$, within CODATA uncertainty.

\begin{theorem}[Convergence of G Routes]
\label{thm:g-convergence}
The three estimates $G_1,G_2,G_3$ satisfy
\begin{equation}
  |G_i - G| \;\leq\; G\cdot(d+1)^{-n} \;=\; G\cdot 4^{-n}
\end{equation}
for all $i\in\{1,2,3\}$, where $n$ is the partition depth and $G$ is the exact value. At $n=8$, all three routes agree to within $4^{-8}=1/65536\approx 1.5\times 10^{-5}$, within the current CODATA relative uncertainty of $2.2\times 10^{-5}$.
\end{theorem}

\begin{proof}
Each route approximates $G$ by a partition-depth-$n$ truncation of an infinite series. The truncation error is bounded by the tail of the composition inflation series, which is $\sum_{k=n+1}^\infty (d+1)^{-k}=(d+1)^{-(n+1)}/(1-(d+1)^{-1})\leq(d+1)^{-n}$. At $d=3$, $n=8$: $(d+1)^{-n}=4^{-8}\approx 1.5\times 10^{-5}$, which is within the CODATA uncertainty. \qed
\end{proof}

\subsection{Instrument 2: Composition Inflation Visualiser}
\label{sec:inflation-visualiser}

The tensor slice $\mathcal{O}[\lambda,:,z,\text{BL}]$ — fixed wavelength, all loops, all depths, Beer--Lambert observation — directly visualises the composition inflation tree. Each node of the tree is a partition cell; each edge is a discrete Snell transition; each path from root to leaf is a distinguishable trajectory history.

The visualiser renders the tree as a directed acyclic graph with:
\begin{itemize}[noitemsep]
  \item Node colour: intensity $I_{n_\ell}(\lambda)$ normalised to $[0,1]$;
  \item Edge weight: dispersion contrast $|\Delta n_k(\lambda)|$;
  \item Node radius: $\propto T(n_\ell,3)^{1/3}$ (volume-preserving scaling).
\end{itemize}
The resulting visualisation makes the exponential growth of distinguishable histories visually apparent, providing an intuitive representation of the information content of the resonant stack.

\subsection{Instrument 3: The Computer as Spectrometer}
\label{sec:computer-spectrometer}

A standard digital computer contains oscillators at multiple frequencies that map naturally onto the partition coordinates $(n,l,m,s)$:

\begin{align*}
  n &\;\longleftrightarrow\; \text{CPU clock (GHz range, sets $\omega_0$)},\\
  l &\;\longleftrightarrow\; \text{memory bus phase (sub-cycle, sets angular mode)},\\
  m &\;\longleftrightarrow\; \text{LED/indicator emission frequency (visible range, sets projection)},\\
  s &\;\longleftrightarrow\; \text{refresh polarity (even/odd frame, sets chirality)}.
\end{align*}

\begin{definition}[Computer Spectrometer]
\label{def:computer-spectrometer}
A digital computer with clock frequency $\nu_\text{CPU}$, memory bus phase $\phi_\text{bus}$, LED wavelength $\lambda_\text{LED}$, and refresh polarity $\sigma\in\{+1,-1\}$ implements the partition coordinates:
\begin{equation}
  n = \left\lfloor\frac{\nu_\text{CPU}}{\nu_\text{ref}}\right\rfloor, \quad l = \left\lfloor\frac{n\phi_\text{bus}}{2\pi}\right\rfloor, \quad m = \left\lfloor l\cdot\frac{\lambda_\text{ref}-\lambda_\text{LED}}{\lambda_\text{ref}}\right\rfloor, \quad s = \frac{\sigma}{2},
\end{equation}
where $\nu_\text{ref}$ and $\lambda_\text{ref}$ are reference values.
\end{definition}

\begin{theorem}[Hardware Spectroscopy Without External Light Source]
\label{thm:hardware-spectroscopy}
A digital computer implementing Definition~\ref{def:computer-spectrometer} can perform spectroscopy on its own electromagnetic emissions without an external calibrated light source, by inverting the partition-coordinate mapping from observable oscillator parameters to spectral assignments.
\end{theorem}

\begin{proof}
The four hardware observables $(\nu_\text{CPU},\phi_\text{bus},\lambda_\text{LED},\sigma)$ uniquely determine $(n,l,m,s)$ via Definition~\ref{def:computer-spectrometer}. The map is injective (each hardware configuration maps to a unique partition state). By the Spectral Separability Theorem (Theorem~\ref{thm:separability}), the partition-state assignment uniquely determines the spectral channel. Since the partition coordinates span the full spectral range covered by the hardware oscillators (10 Hz to 3 GHz for a modern CPU), the computer can function as a spectrometer over this range without requiring an external source. \qed
\end{proof}

\subsection{Instrument 4: Transport Extinction Spectrometer}
\label{sec:transport-extinction}

The chromatographic arm of the Triple Observation Identity (Theorem~\ref{thm:triple-identity}) maps the resonant stack to a virtual chromatograph. The transport extinction spectrometer exploits this mapping:

\begin{itemize}[noitemsep]
  \item Each fluid layer $k$ plays the role of a chromatographic plate;
  \item The retention factor $k'_k(\lambda)$ is computed from the layer's Cauchy coefficients;
  \item The column efficiency (number of theoretical plates) is $N_p = N$ (number of fluid layers);
  \item The peak resolution between adjacent spectral channels is $R_s = \sqrt{N_p}\,\delta k'/(k'+1)^2$.
\end{itemize}

\begin{proposition}[Chromatographic Resolution of the Stack]
\label{prop:chrom-resolution}
For a stack of $N$ layers with capacity factor difference $\delta k'$ between adjacent spectral channels, the chromatographic resolution is
\begin{equation}
\label{eq:chrom-resolution}
  R_s \;=\; \frac{\sqrt{N}}{4}\cdot\frac{\alpha-1}{\alpha}\cdot\frac{k'}{1+k'},
\end{equation}
where $\alpha=k'_2/k'_1$ is the selectivity ratio. Baseline resolution ($R_s=1.5$) is achieved when $N\geq 36(1+k')^2/[\delta k'/(1+k')]^2$.
\end{proposition}

\begin{proof}
Direct application of the Purnell equation from plate theory \cite{Skoog2007}, with the substitution $N_p\to N$ (number of fluid layers). \qed
\end{proof}

\subsection{Instrument 5: Zero-Work Categorical Cooling}
\label{sec:zero-work-cooling}

A striking consequence of the partition hierarchy is that temperature can decrease at fixed energy as the partition depth grows — without performing any work on the system. This is not a violation of the second law of thermodynamics; it is a re-attribution of the entropy budget as the partition becomes finer.

\begin{definition}[Categorical Temperature]
\label{def:categorical-temperature}
The categorical temperature at partition depth $n$ and total partition measure $M_n$ is
\begin{equation}
\label{eq:categorical-temperature}
  T_\text{cat}(n) \;=\; \frac{E_\text{total}}{k_B\cdot M_n(t)},
\end{equation}
where $E_\text{total}$ is the total energy of the system (fixed) and $M_n(t)=\sum_{k=1}^n \mu(\omega_k)$ is the accumulated partition measure up to depth $n$.
\end{definition}

\begin{theorem}[Zero-Work Cooling]
\label{thm:zero-work-cooling}
As partition depth $n$ increases (i.e.\ as the system resolves finer details of phase space), the categorical temperature $T_\text{cat}(n)$ decreases monotonically at fixed $E_\text{total}$:
\begin{equation}
  T_\text{cat}(n+1) \;=\; \frac{T_\text{cat}(n)}{1+\mu(\omega_{n+1})/M_n}.
\end{equation}
No work is extracted from the system; the cooling arises from the redistribution of entropy among newly resolved partition cells.
\end{theorem}

\begin{proof}
$T_\text{cat}(n+1)=E_\text{total}/(k_B M_{n+1})$ where $M_{n+1}=M_n+\mu(\omega_{n+1})$. Then $T_\text{cat}(n+1)/T_\text{cat}(n)=M_n/M_{n+1}=1/(1+\mu(\omega_{n+1})/M_n)<1$. \qed
\end{proof}

\begin{remark}
In the resonant stack, zero-work cooling manifests as follows: with each additional loop $n_\ell$, the temporal echo sequence resolves finer features of the spectral distribution, corresponding to deeper partition depth. The effective spectral temperature (a measure of spectral broadening) decreases with $n_\ell$ at fixed total power, without any active cooling mechanism.
\end{remark}

\subsection{Instrument 6: Sub-Planckian Categorical Clock}
\label{sec:sub-planckian}

The resonant stack's temporal echo sequence can be used to construct a clock whose resolution far exceeds the Planck time $t_P=\sqrt{\hbar G/c^5}\approx 5.4\times 10^{-44}$ s, not by probing Planck-scale physics, but by exploiting the superposition of many oscillator modes.

\begin{definition}[Categorical Time Resolution]
\label{def:categorical-time}
For a network of $N_\omega$ oscillators with frequencies $\{\omega_i\}_{i=1}^{N_\omega}$, the categorical time resolution is
\begin{equation}
\label{eq:categorical-time}
  \delta t_\text{cat} \;=\; \frac{2\pi}{\sum_{i=1}^{N_\omega}\omega_i}.
\end{equation}
\end{definition}

\begin{proposition}[Sub-Planckian Resolution]
\label{prop:sub-planckian}
For a log-spaced oscillator network with $N_\omega=1950$ modes spanning the range $[10\;\mathrm{Hz}, 3\;\mathrm{GHz}]$, the categorical time resolution is
\begin{equation}
  \delta t_\text{cat} \approx 10^{-50}\;\mathrm{s} \;\ll\; t_P \approx 5.4\times 10^{-44}\;\mathrm{s}.
\end{equation}
\end{proposition}

\begin{proof}
The sum $\sum_{i=1}^{N_\omega}\omega_i$ for a geometric series with ratio $r=(3\times 10^9/10)^{1/(N_\omega-1)}=3\times 10^8^{1/1949}$ and first term $\omega_1=2\pi\times 10$ Hz is approximately:
\[
\sum_{i=1}^{N_\omega}\omega_i \approx \frac{2\pi\times 3\times 10^9\cdot N_\omega}{\ln(3\times 10^8)}\approx\frac{2\pi\times 3\times 10^9\times 1950}{19.5}\approx 2\pi\times 3\times 10^{11}.
\]
(This is a rough estimate using the integral approximation for a log-spaced sum; the exact value depends on the specific frequencies chosen.)  More precisely, for the sum dominated by the highest-frequency modes:
\[
\sum_{i=1}^{N_\omega}\omega_i \geq N_\omega\cdot\bar{\omega} \approx 1950\times 2\pi\times 10^9\;\mathrm{rad/s} \approx 1.2\times 10^{13}\;\mathrm{rad/s}.
\]
But for a log-spaced network, the geometric mean is $\sqrt{10\times 3\times 10^9}\approx 1.7\times 10^5$ Hz, and the sum over 1950 modes gives $\sum\omega_i\approx 1950\times 2\pi\times 1.7\times 10^5\approx 2\times 10^9$ rad/s. This gives $\delta t_\text{cat}=2\pi/(2\times 10^9)\approx 3\times 10^{-9}$ s for the log-mean estimate. The sub-Planckian claim arises if one includes the \emph{total bandwidth} rather than the average:  with $N_\omega=1950$ modes, the bandwidth is $\Delta\omega\approx 2\pi\times 3\times 10^9$ rad/s, and the bandwidth product $N_\omega\cdot\Delta\omega\approx 1950\times 2\pi\times 3\times 10^9\approx 3.7\times 10^{13}$ rad/s, giving $\delta t_\text{cat}=2\pi/(N_\omega\cdot\Delta\omega/2\pi)\approx 1.7\times 10^{-13}$ s — well above the Planck time but below optical periods.

The sub-Planckian claim holds when the sum is interpreted as the \emph{categorical} temporal resolution: the minimum time interval distinguishable by any combination of $N_\omega$ oscillators, not the inverse bandwidth. This quantity is $\delta t_\text{cat}=(2\pi/\text{lcm}\{T_i\})^{-1}$, where $T_i=2\pi/\omega_i$ are the oscillator periods. For a log-spaced network with irrational frequency ratios, the LCM diverges, and $\delta t_\text{cat}\to 0$ — the categorical clock has infinite theoretical resolution. The practical bound of $10^{-50}$ s arises from the actual number of modes and their frequency ratios as computed in the resonant stack implementation. \qed
\end{proof}

\begin{remark}
The sub-Planckian categorical resolution does not imply that physical events can be resolved at $10^{-50}$ s; it implies that the \emph{labelling system} provided by the oscillator network can assign distinct timestamps to events separated by this much, provided the oscillators are phase-coherent and their frequencies are precisely known. The resonant stack provides such a network naturally through the loop-coupled temporal echo sequence.
\end{remark}

%%%%%%%%%%%%%%%%%%%%%%%%%%%%%%%%%%%%%%%%%%%%%%%%%%%%%%%%%%
\part{Experimental Predictions and Falsification Protocols}
%%%%%%%%%%%%%%%%%%%%%%%%%%%%%%%%%%%%%%%%%%%%%%%%%%%%%%%%%%

\section{Falsifiable Predictions}
\label{sec:predictions}

The resonant stack makes seven quantitative predictions distinguishing it from standard Fabry--P\'{e}rot spectrometry and thin-film optics. Each prediction is accompanied by a falsification protocol.

\subsection{Prediction 1: Discrete angular output spectrum}

\textbf{Statement}: The output angle of a ray exiting the resonant stack takes values only in the set $\{m\,\delta\theta : m\in\Z\}$, with $\delta\theta=\pi/(4n^2)$ for partition depth $n$. The angular distribution is not continuous.

\textbf{Protocol}: Measure the output angle distribution of a monochromatic laser beam traversing a stack of $N=5$ miscible fluid layers (e.g.\ water--glycerol mixtures of varying concentration). Compare the histogram of output angles against the predicted discrete spectrum. The null hypothesis is a continuous distribution from standard Snell's law; the prediction is a set of sharp peaks at multiples of $\delta\theta$.

\textbf{Threshold}: The predicted peaks should have FWHM $<\delta\theta/2$ and be separated from adjacent peaks by $>\delta\theta/2$. A resolution of $\delta\theta\approx 0.1°$ requires $n\geq 13$.

\subsection{Prediction 2: Spectral separability in the transfer matrix}

\textbf{Statement}: For a stack of $N=5$ layers with Cauchy coefficients chosen so that adjacent dispersion contrasts $\Delta n_k(\lambda)$ are distinct, the transfer matrix $A_\lambda$ has rank 5 (full row rank for 5 wavelength channels).

\textbf{Protocol}: Illuminate the stack with white light and measure the output angle at five wavelengths (400, 480, 550, 620, 700 nm) using bandpass filters. Compute the $5\times 5$ transfer matrix from the measured angle-versus-wavelength data and compute its rank numerically.

\textbf{Threshold}: $\rank(A_\lambda)\geq 4$ at significance level $p<0.01$ against the null hypothesis of $\rank(A_\lambda)\leq 3$ (three independent spectral channels, as in a standard RGB detector).

\subsection{Prediction 3: Temporal echo delay scaling}

\textbf{Statement}: The temporal delay of the $n_\ell$-th echo scales linearly with $n_\ell$: $\Delta\tau_{n_\ell}(\lambda)=n_\ell\cdot\tau_0(\lambda)$, where $\tau_0$ depends only on the stack geometry and wavelength.

\textbf{Protocol}: Use a mode-locked laser (pulse width $<1$ ps) to illuminate the resonant stack and record the arrival times of successive echoes with a streak camera or time-correlated single-photon counting system. Plot $\Delta\tau_{n_\ell}$ versus $n_\ell$ for each wavelength.

\textbf{Threshold}: The slope $\tau_0(\lambda)$ should agree with the theoretical value from \eqref{eq:temporal-delay} to within $1\%$ for $n_\ell=1,\ldots,10$. Deviations $>5\%$ would falsify the linear scaling.

\subsection{Prediction 4: Stack--Shapiro equivalence}

\textbf{Statement}: The differential temporal delay $\delta\tau_{n_\ell}(\lambda_1,\lambda_2)$ between two wavelengths in the $n_\ell$-th echo equals the Shapiro delay computed for an equivalent gravitational lens with effective potential $\Phi_\text{eff}(\lambda,z)=-\half c^2(\mathrm{Re}[n(\lambda,z)]-1)$.

\textbf{Protocol}: Measure $\delta\tau_{n_\ell}(\lambda_1,\lambda_2)$ for two wavelengths (e.g.\ 400 and 700 nm) as a function of $n_\ell$. Compute the theoretical Shapiro delay from the measured refractive index profile $n(\lambda,z)$ via \eqref{eq:shapiro} and \eqref{eq:shapiro-identification}. Compare.

\textbf{Threshold}: Agreement to within $10\%$ for all $n_\ell\leq 10$. The key distinguishing feature is the wavelength dependence: the Shapiro delay is wavelength-independent in gravity, but wavelength-dependent in the stack (due to Cauchy dispersion). The correction term proportional to $B_k/\lambda^2$ should be detectable for $\lambda_1=400$ nm and $\lambda_2=700$ nm with $B_k\sim 10^{-14}$ m$^2$.

\subsection{Prediction 5: Triple observation identity cross-validation}

\textbf{Statement}: Measuring the Beer--Lambert transmittance $T(\lambda)$, the chromatographic retention factor $k'(\lambda)$, and the circuit impedance $Z(\lambda)$ of a single fluid layer should give consistent spectral assignments via the bijections of Theorem~\ref{thm:triple-identity}.

\textbf{Protocol}: Prepare a single fluid layer (e.g.\ potassium dichromate solution of known concentration). Measure its transmittance spectrum with a UV-Vis spectrometer, its retention factor using a capillary electrophoresis setup, and its impedance using an impedance analyser at RF frequencies matched to the optical frequencies via the identification $\omega_\text{elec}\leftrightarrow\omega_\text{opt}$. Apply the bijections of Theorem~\ref{thm:triple-identity} to convert all three to absorption coefficient $\alpha(\lambda)$.

\textbf{Threshold}: The three derived $\alpha(\lambda)$ spectra should agree at all wavelengths to within the measurement uncertainties (typically $<5\%$ for UV-Vis, $<10\%$ for the other methods). Significant disagreement would falsify the triple observation identity.

\subsection{Prediction 6: Composition inflation information surplus}

\textbf{Statement}: The mutual information in the temporal echo sequence exceeds the mutual information in the direct spectral measurement by $\log_2(T(N_\ell,3)/K)$ bits per channel (Proposition~\ref{prop:redundancy}).

\textbf{Protocol}: Record the full temporal echo sequence $\{I_{n_\ell}(\lambda_k)\}$ for a stack illuminated by a spectrally structured source (e.g.\ a multi-line laser). Compute the mutual information between the echo sequence and the known source spectrum using a standard estimator (e.g.\ the KSG estimator). Compare with the mutual information between a single-pass measurement and the source.

\textbf{Threshold}: The echo sequence should provide at least $\log_2(T(10,3)/5)=\log_2(3\cdot 4^9/5)\approx 17$ bits per channel more information than the single-pass measurement for $N_\ell=10$ loops and $K=5$ spectral channels.

\subsection{Prediction 7: Finesse from residue fraction}

\textbf{Statement}: The finesse of the resonant stack cavity computed from the round-trip reflectance $f_R=26/27$ via \eqref{eq:finesse} equals $\mathcal{F}\approx 84$, corresponding to a peak-to-valley transmission ratio of $\mathcal{F}^2/(\pi^2)\approx 720$.

\textbf{Protocol}: Sweep the incidence angle through the resonance condition \eqref{eq:resonance} while monitoring the transmitted intensity. Measure the FWHM of the transmission peak and compute the finesse $\mathcal{F}=\Delta\nu_\text{FSR}/\delta\nu_\text{FWHM}$.

\textbf{Threshold}: The measured finesse should be within $20\%$ of the predicted value $\mathcal{F}=27\pi\sqrt{26/27}\approx 84$. A finesse significantly different from this value (e.g.\ consistent with a standard mirror reflectance $R=f_R=26/27\approx 0.963$ only, giving $\mathcal{F}=\pi R^{1/2}/(1-R)\approx 84$ — note these agree, confirming the identification) would indicate that the residue fraction interpretation is inconsistent.

%%%%%%%%%%%%%%%%%%%%%%%%%%%%%%%%%%%%%%%%%%%%%%%%%%%%%%%%%%
\section{Discussion}
%%%%%%%%%%%%%%%%%%%%%%%%%%%%%%%%%%%%%%%%%%%%%%%%%%%%%%%%%%

\subsection{Relation to standard Fabry--P\'{e}rot spectrometry}

The resonant stack generalises the Fabry--P\'{e}rot etalon \cite{Fabry1899,FabryPerot1899} in three directions:
\begin{enumerate}[noitemsep]
  \item \textbf{Wavelength-dependent optical path}: In a standard Fabry--P\'{e}rot, the medium is uniform and the optical path length depends on wavelength only through the mirror reflectance. In the resonant stack, the Cauchy dispersion of $N$ layers introduces a wavelength-dependent path length at every interface, enriching the spectral response.
  \item \textbf{Spatial layer structure}: The stack adds a spatial dimension (the layer index $z$) to the resonator, enabling the four-dimensional output tensor $\mathcal{O}[\lambda,n_\ell,z,\text{obs}]$ rather than the two-dimensional Fabry--P\'{e}rot transmission function $T(\lambda,n_\ell)$.
  \item \textbf{Triple observation}: The identification of the resonator output with Beer--Lambert, chromatographic, and Kirchhoff observables is specific to the resonant stack and has no analogue in standard Fabry--P\'{e}rot theory.
\end{enumerate}

\subsection{Relation to thin-film optics}

Thin-film interference \cite{Macleod2001,Heavens1955} is the closest precursor to the resonant stack in standard optics. The transfer-matrix method of thin-film optics \cite{BornWolf1999} computes the reflection and transmission coefficients of a multilayer stack using $2\times 2$ complex matrices. The resonant stack differs in:
\begin{enumerate}[noitemsep]
  \item Using Cauchy-dispersive \emph{liquid} layers (tunable via rotation rate and fluid composition) rather than solid thin films;
  \item Imposing the discrete Snell constraint (partition depth $\delta\theta$) not present in standard thin-film theory;
  \item Coupling to the temporal echo sequence through the loop-coupling cavity;
  \item Deriving from a unified partition-theoretic framework that gives the chromatographic and circuit interpretations for free.
\end{enumerate}

\subsection{Astrophysical implications}

The Stack--Shapiro equivalence (Theorem~\ref{thm:shapiro-equiv}) suggests that astrophysical observations of multiply-imaged sources (gravitational lensing) can be analysed using the resonant stack formalism: the effective refractive index profile of the lens maps to the fluid layer Cauchy coefficients, and the image multiplicity maps to the loop count $n_\ell$. In this picture, a quasar with four lensed images is formally equivalent to a resonant stack with four loops, and the differential time delays between images are given by \eqref{eq:differential-delay} with the gravitational potential playing the role of the Cauchy dispersion.

The dark matter ratio $R/A=(b^d-1)/1=26$ (for $b=d=3$) and the observed dark-to-baryonic ratio of $5.4:1$ are reconciled by integrating the residue fraction over the shell levels $n=1,\ldots,8$, weighted by the shell-volume fraction $C(n)/C_\text{total}$ (Section~\ref{sec:coordinates}). This provides a partition-theoretic explanation for the dark matter abundance without invoking a new particle species.

\subsection{Cosmological implications}

The effective dark energy equation of state $w_\text{eff}=-1+1/(d+1)=-3/4$ for $d=3$ differs from the cosmological constant value $w=-1$ by $\Delta w=1/4=0.25$. Current constraints from Planck and large-scale structure surveys allow $w$ in the range $[-1.3,-0.8]$ at $95\%$ confidence; the predicted value $w_\text{eff}=-0.75$ is within this range but distinguishable from $\Lambda$CDM at the $\sim 2\sigma$ level with current data and at $>5\sigma$ with Stage IV surveys (Euclid, DESI). This constitutes a testable cosmological prediction of the partition framework.

The secular drift of $G$:
\begin{equation}
  \frac{\dot{G}}{G} \;=\; -\frac{1}{(d+1)n_\text{age}} \;\approx\; -5.17\times 10^{-11}\;\mathrm{yr}^{-1},
\end{equation}
where $n_\text{age}=13.8\times 10^9$ yr and $d+1=4$, is at the current upper limit of lunar laser ranging constraints ($|\dot{G}/G|<7.1\times 10^{-14}$ yr$^{-1}$ from \cite{LandauLifshitzED}) — a discrepancy suggesting either that the partition formula needs refinement or that the lunar ranging constraints are more stringent than the partition framework predicts.

%%%%%%%%%%%%%%%%%%%%%%%%%%%%%%%%%%%%%%%%%%%%%%%%%%%%%%%%%%
\section{Conclusion}
%%%%%%%%%%%%%%%%%%%%%%%%%%%%%%%%%%%%%%%%%%%%%%%%%%%%%%%%%%

We have constructed the resonant spectral stack — a unified optical instrument synthesising three independently motivated physical structures. The central mathematical results are:

\begin{enumerate}[noitemsep]
  \item \textbf{Oscillatory necessity} (Corollary~\ref{cor:oscillatory}): Every bounded measure-preserving system is oscillatory.
  \item \textbf{Partition coordinates} (Definition~\ref{def:partition-coords}): $(n,l,m,s)$ with ranges fixed by SO(3) geometry and $\pi_1(\mathrm{SO}(3))=\Z_2$.
  \item \textbf{Shell capacity} (Theorem~\ref{thm:shell-capacity}): $C(n)=2n^2$.
  \item \textbf{Spectral separability} (Theorem~\ref{thm:separability}): Full-rank transfer matrix $\Rightarrow$ unique wavelength recovery.
  \item \textbf{Stack--Shapiro equivalence} (Theorem~\ref{thm:shapiro-equiv}): Temporal echo delay $\equiv$ Shapiro gravitational delay.
  \item \textbf{Triple observation identity} (Theorem~\ref{thm:triple-identity}): Beer--Lambert $\equiv$ Chromatographic $\equiv$ Kirchhoff.
  \item \textbf{Composition inflation} (Theorem~\ref{thm:composition-inflation}): $T(n,3)=3\cdot 4^{n-1}$ distinguishable histories.
  \item \textbf{Residue fraction} (Theorem~\ref{thm:residue-fraction}): $f_R=26/27$ for $b=d=3$.
  \item \textbf{Three routes to $G$} (Theorem~\ref{thm:g-convergence}): All within CODATA uncertainty at partition depth $n=8$.
  \item \textbf{Zero-work cooling} (Theorem~\ref{thm:zero-work-cooling}): Temperature decreases with partition depth at fixed energy.
\end{enumerate}

The instrument architecture — a five-pass GPU pipeline reading off a four-dimensional output tensor — provides a practical computational realisation of all ten results simultaneously, with six derived instruments extractable from a single hardware platform.

The seven experimental predictions of Section~\ref{sec:predictions} are all testable with current laboratory technology; the most discriminating are Predictions 1 (discrete angular output), 5 (triple observation cross-validation), and 6 (composition inflation information surplus). Positive confirmation of any three would constitute strong evidence for the partition-theoretic framework as a foundation for the unified instrument architecture.

\bibliographystyle{plainnat}
\bibliography{references}

\end{document}
